\documentclass[12pt,twoside]{book}
\usepackage[a4paper,width=150mm,top=25mm,bottom=20mm,headheight=14.5pt]{geometry}
\usepackage[natbib,style=apa,backend=biber,backref=true]{biblatex}
\bibliography{thesis.bib}
\usepackage{fancyhdr}
\pagestyle{fancy}
\usepackage{algorithm}% http://ctan.org/pkg/algorithm
\usepackage{algpseudocode}% http://ctan.org/pkg/algorithmicx
\usepackage{amsmath,amsfonts,amssymb}
\usepackage{caption}
\usepackage{subcaption}
\usepackage{CJKutf8}
\usepackage[T1]{fontenc}
\usepackage[acronym,symbols,nogroupskip,nonumberlist,toc]{glossaries-extra}
\usepackage{graphicx}
\graphicspath{ {figures/} }
% \DeclareGraphicsExtensions{.png,.pdf,.eps}  % use during development
\DeclareGraphicsExtensions{.eps,.pdf,.png}  % use for final production
\usepackage[bookmarks]{hyperref}
\usepackage[utf8]{inputenc}
\usepackage{lscape}
\usepackage{mathtools}
\usepackage{siunitx}
\usepackage{tabularx}
\usepackage{textcomp}
\usepackage{titlesec}

% SIrange using em-dash
\sisetup{range-phrase=--}
\sisetup{range-units=single}

% subsubsubsections paragraphs
\setcounter{secnumdepth}{3}
\titleformat{\paragraph}
{\normalfont\normalsize\bfseries}{\theparagraph}{1em}{}
\titlespacing*{\paragraph}
{0pt}{3.25ex plus 1ex minus .2ex}{1.5ex plus .2ex}

% glossary items
\makeglossaries[\acronymtype]

\glsxtrnewsymbol[description={hydropotential}]{phi}{\ensuremath{\phi}}
\glsxtrnewsymbol[description={time}]{t}{\ensuremath{t}}
\newacronym{ConvNets}{ConvNets}{Convolutional Neural Networks}
\newacronym{DBSCAN}{DBSCAN}{Density-based Spatial Clustering of Applications with Noise}
\newacronym{DEM}{DEM}{Digital Elevation Model}
\newacronym{ESRGAN}{ESRGAN}{Enhanced Super-Resolution Generative Adversarial Network}
\newacronym{ICESat-2}{ICESat-2}{
  Ice, Cloud, and land Elevation Satellite 2. A near polar-orbit laser altimeter satellite that measures ice surface elevation.
}
\newacronym{IPCC}{IPCC}{Intergovernmental Panel on Climate Change}
\newacronym{ISSM}{ISSM}{Ice-sheet and Sea-level System Model}
\newacronym{LeakyReLU}{LeakyReLU}{Leaky Rectified Linear Unit}
\newacronym{REMA}{REMA}{
  Reference Elevation Model of Antarctica. A high resolution, time-stamped Digital Surface Elevation Model of Antarctica created from stereophotogrammetry.
}
\newacronym{RES}{RES}{Radio-echo sounding}
\newacronym{SD}{SD}{Standard deviation}
\newacronym{SLE}{SLE}{Sea level equivalent}
%Not a Number (NaN)

\DeclareSIUnit\year{yr}

\title{
{The subglacial landscape and hydrology of Antarctica as mapped from space}\\
{\large Antarctic Research Centre, Victoria University of Wellington}\\
{\includegraphics[scale=0.5]{0_arc_logo.jpg}}
}
\author{Wei Ji Leong}
\date{\today}

\begin{document}

\begin{titlepage}
  \begin{center}
    \vspace*{1cm}

    \Huge
    \textbf{Deep Learning in the Cryosphere}

    \vspace{0.5cm}
    \LARGE
    The subglacial landscape and hydrology of Antarctica as mapped from space

    \vspace{1.5cm}

    \textbf{Wei Ji Leong}

    \vfill

    A thesis presented for the degree of\\
    Doctor of Philosophy

    \vspace{0.8cm}

    \includegraphics[width=0.4\textwidth]{arc_logo}

    \Large
    Antarctic Research Centre\\
    Victoria University of Wellington\\
    New Zealand\\
    \today

  \end{center}
\end{titlepage}


\chapter*{Abstract}

To narrow uncertainties in the Antarctic ice sheet's contribution to sea level rise, we present a collection of novel machine learning and automated satellite remote sensing methods which use ice surface observations to infer the subglacial nature of Antarctica.
A super-resolution deep neural network called DeepBedMap was designed and trained to produce a high-resolution (\SI{250}{\metre}) bed elevation model of Antarctica called DeepBedMap\_DEM that preserves bed roughness details useful for catchment- to continent-scale ice sheet modelling.
This DeepBedMap\_DEM is compared with a smoother, medium-resolution (\SI{500}{\metre}) BedMachine topography in a basal inversion experiment over Pine Island Glacier, with results motivating more research into the interacting roles of subglacial hydrology which influences skin drag and high resolution bed topographies which influences form drag.
Active subglacial lakes in Antarctica were mapped using an unsupervised density-based classification method on ICESat-2 point cloud data from 2018-2020, yielding 195 active subglacial lakes, including 36 new lakes in the 86--88°S area not detected by the previous ICESat (2003-2009) mission.
This thesis showcases both the rich diversity in subglacial landscapes and the dynamic nature of subglacial hydrology in Antarctica, forming a foundation enabling the accurate modelling of overland ice flow in critical regions of the vulnerable West Antarctic Ice Sheet.

\section*{Plain language summary}

Antarctica has a lot of ice, but we're unsure how fast ice can slide into the sea and cause water to go up in beaches around the world.
So we teach computers to solve hard math problems that tell us how fast sea water might go up.
These computers are fed with lots of pictures taken from cameras up in the sky and space.
Ice sits on top of rock in Antarctica, and with practice, the computers get pretty good at telling us how high and bumpy the rock is.
The rock under the ice appears quite bumpy, and ice probably doesn't like sliding over bumpy rocks since it's rough.
Sometimes though, ice may not mind sliding over rough bits of rock if the rock moves along with it, or if water gets in between the rock and ice to makes things slippery, but we ask our smart computers to be sure.
There are also lasers from space shooting down at earth and bouncing back to tell us how ice in Antarctica is going up or down.
Once in a while, they tell us that ice in parts of Antarctica moved up or down a bit too fast.
Smart people think these are lakes hiding under the ice, filling up with water or draining, and we found many of these lakes over Antarctica, especially in an area called Whillans Ice Stream on the Siple Coast.
We hope that the computers can keep learning faster because there's a lot of pictures showing ice moving pretty fast, and it doesn't look like there's much time before a big chunk of ice might break away in Antarctica and flood beaches around the world.

\chapter*{Dedication}
To Penny, who led me, so profoundly, to be

\chapter*{Acknowledgements}

\begin{CJK*}{UTF8}{gbsn}
每一个成功男人的背后，都有一个伟大的女人。
Every step along this journey has not been without the support of a great community behind me.
My Chinese name is 梁伟吉，each character signifying a part of who I am, my duty in this world, and the enormous privillege that has allowed me to make it to this day.

Family 粱 (Leong) comes first.
My dearest partner Penny Qiu Yue 邱乐, you have went above and beyond what a normal long distance relationship entails, keeping me company daily through good times and bad, being patient, trusting and more than I could ever express in my supposedly native yet broken Malaysian/Bruneian Chinese.
Above all, thank you for sharing the most precious asset there is - 时间 - time - which for that, I am forever grateful.
To my younger sister Wei Qian 伟倩 and brother Wei Tong 伟桐，growing up with you two in Brunei has taught me the importance of sharing and looking after each other, so thank you for your understanding, for all the catch up calls, and simply for being.
For Dad and Mum，梁耀明 and 陈碹丝，it takes a million to raise one child, let alone three, and I have lost count of the number of side jobs you've had over the years, especially the struggles in recent times, both of you deserve a special thank you and a long holiday.

To my local family in Wellington, Jaclyn and David Chai, for the nourishment you have provided over the years in the form of yummy Malaysian style vegetarian food.
It is unbelievable how fast time flies, seeing your two kids ZH and Yi Ling grow tall and mature, you are all my second family and will always be remembered.

Greatness 伟 (Wei) of teachers is next.
To my great supervisor Huw Horgan (and his amazing partner Ruschle), I sometimes wonder how a person from the tropics without a trace of snow got into studying Antarctic ice, and obviously it was because of you, from that time in 2014 studying proglacial lakes at Mount Cook national park, to my 2015 Honours project on mapping potential subglacial estuaries in Antarctica, to this very PhD thesis 2017-2021.
A very big thank you for allowing me to do my work remotely using satellites, sending me randomly on exchanges to Alaska and Austria to realize how big the world truly is, for giving me the space to develop independently, while providing top notch Yoda-like advice that takes me weeks and months to figure out.

Thank you to all my Geography teachers and mentors over the years.
In particular, to my secondary school teacher Veronica Chan at Chung Hwa Middle School, for your great style of Geography teaching with funny stories and sketches, it was weird that you didn't recommend me taking Geography in high school, but here I am now.
% To my high school teacher
To my university lecturer Mairéad de Róiste, for putting me on this journey into Geographic Information Systems (GIS), setting good examples on how to tutor the next generation of students, and also for drilling in Cartogrphic best practices.
To Karl Baker at e-Spatial, for passing on your guru level of stacking and mosaicking aerial imagery, arranging map layers to perfection, and for being a role model on how to be a respectable GIS professional.

Luck 吉 (Ji) is the last but not least.
I feel really fortunate to be able to stand on the shoulders of giants, and would like to acknowledge at least some of the many people I've met along this journey.

\paragraph{Open source}

To Leonardo Uieda, DongDong Tian 田冬冬, Liam Toney and the rest of the PyGMT team, and also to Paul Wessel for his enduring work on the Generic Mapping Tools (GMT) package over many decades, it is a pleasure to have such a great tool for producing beautiful maps。
To all the folks maintaining the Python and Pangeo ecosystem of geospatial packages, from Xarray to RAPIDS AI and every imaginable library in between, thank you for all the new features and bug fixes, keep up the good work!

\paragraph{Europe}

To Poonam Vishwas, KaiKai Xue 薛凯, Yawen Tan 谭雅文, and all the friends I met in the 2019 Erasmus+ exchange at ZGIS, University of Salzburg, Austria, it was an awesome experience having so many like-minded people researching GIS together in Europe, an experience I will always treasure.
To Dr. Hermann Klug, thank you for warmly inviting me to Salzburg and sorting out all manners of logistics in German, and to Dr. Dirk Tiede for introducing me to Object Based Image Analysis and providing a conceptual framework for understanding the place of Convolutional Neural Networks in the wider field supervised classification is GIS.

\paragraph{United States}

To Anushilan Acharya, Indraneel Kasmalkar, Whyjay Zheng, and all my friends in the 2018 McCarthy International Summer School in Glaciology, it was a great pleasure meeting you all, and I look forward to seeing some of you again at some point!
To Dr. Regine Hock, thank you for inviting me to Alaska and giving me a once in a lifetime experience of stepping on a big glacier, and to Dr. Kelly Brunt for some great lessons on the ICESat-2 laser altimetry satellite!
To the International Glaciological Society (IGS), thank you for the generous early career travel grant and the chance to present a talk at the 2019 IGS Five Decades of Radioglaciology Symposium at Stanford University.

\paragraph{New Zealand}

To Arran Whiteford, Alanna Alevropoulos-Borrill, Francesca Baldacchino, my CO522 office mates at the Antarctic Research Centre, for all the yarns about ice and other topics, I couldn't have asked for a nicer group of friendly folks to do my PhD with!
To Dao Polsiri and Michelle Dow, thank you for sorting out all sorts of paperwork from flight ticket changes to funding for trivial things, and also to Aleksandr Beliaev, for troubleshooting issues with the UNIX servers and putting up with my overuse of storage, RAM, GPU resources and all.

Finally to anyone that I left out, you are not unnoticed, this sentence is for you, and I hope you enjoy reading this thesis just as much as I do.

\end{CJK*}

\tableofcontents
\listoffigures
\listoftables

% \printnoidxglossary[type=symbols,sort=use,style=long,title={List of Symbols}]
\glsaddall
\setglossarystyle{altlist}
\printglossary[type=\acronymtype]

\chapter{Introduction}
% The subglacial landscape and hydrology of Antarctica as mapped from space

\section{Motivation}

What is the rate at which sea level will rise and affect coastal communities around the globe?
Besides greenhouse gas emissions, a large uncertainty in the answer to that question lies in the Antarctic ice sheet's behaviour.
By 2100, Antarctica is projected to contribute as little as \SI{0.03}{\metre} to as much as \SI{0.28}{\metre} sea level equivalent (\gls{SLE}) of water under a high emissions scenario \citep[RCP 8.5, Fig.~\ref{fig:0}f][]{PortnerSummaryPolicymakers2019}.
To narrow the range of this projection, one of the regions we need to understand is Antarctica's bed, the part hidden under the ice.
This thesis synthesizes a collection of satellite remote sensing methods, using what we can see at the surface of the Antarctic continent, to deduce what is happening in the subglacial world below.

\begin{figure}[htbp]
  \includegraphics[width=1.0\textwidth]{chp1_fig_SROCC_SPM1_Final_edit_AA-1.jpg}
  \caption[Cryospheric contribution to sea level change by 2100]{
    Contribution of the cryosphere to sea level change.
    Modelled and observed historical changes in the cryosphere since 1950, and the projected future changes under low (RCP2.6) and high (RCP8.5) greenhouse gas emissions scenarios.
    Figure adapted from SPM.1 in the Summary for Policy Makers, \gls{IPCC} Special Report on the Ocean and Cryosphere in a Changing Climate \citep{PortnerSummaryPolicymakers2019}.
  }
  \label{fig:0}
\end{figure}

\section{Thesis context}

Antarctica is a large continent, $53\times$ the size of New Zealand or $1.8\times$ the size of Australia, and satellite measurements are one of the best ways to cover the whole icescape.
While the first satellite images of Antarctica were taken in 1964 by the Nimbus I satellite \citep{MeierNewestimatesArctic2013}, continuous modern-day observations only started in 1979.
Forty years later, and there is now a wealth of openly available optical images of much of the continent, currently captured by NASA's Landsat 8 \citep[][]{RoyLandsat8Scienceproduct2014} and the European Space Agency (ESA)'s Sentinel 2 \citep{DruschSentinel2ESAOptical2012}.
Changes to Antarctica's mass can be measured at kilometre-scale using the Gravity Recovery and Climate Experiment Follow‐On (GRACE‐FO, 2018-) satellites \citep{LandererExtendingGlobalMass2020}, and at metre-scales using the Ice, Cloud and land Elevation Satellite-2 \citep[ICESat-2, 2018-;][]{MarkusIceCloudland2017}.
Also with radar instruments like Sentinel 1 \citep{AttemaSentinel1missionoverview2009} and Cryosat 2 \citep{WinghamCryoSatmissiondetermine2006}, cloud-free monitoring of the Antarctic surface has never been easier.
Sadly, most of these satellites can only observe the surface of the Antarctic continent.
There are physical constraints on designing an Earth orbiting satellite which can penetrate through both the atmosphere and kilometres of ice \citep{GogineniCubeSatTrainRadar2018,CulbergFirnClutterConstraints2020}.
Arguably more is known about the subglacial landscape of Mars's polar caps \citep{OroseiRadarevidencesubglacial2018,ArnoldModeledSubglacialWater2019,LauroMultiplesubglacialwater2020} than that of our own planet Earth.
Still, it is possible to infer about Antarctica's bed properties by combining satellite observations captured over large surface areas with what little is known about the bed at discrete sites targeted by oversnow observations, and this is the approach this thesis will take.

\subsection{Research Questions}

To better understand the drivers of Antarctic ice flow, we present high spatial resolution ($<= \SI{250}{\metre}$) views of Antarctica's subglacial geography, as inferred from satellite detected ice surface changes.
Novel machine learning and automated remote sensing techniques are applied to infer both the subglacial landscape and subglacial hydrological network of Antarctica.
The 3 main research questions are as follows:

\begin{enumerate}
  \item How can we integrate existing high spatial resolution remote sensing products to boost the resolution of existing bed elevation maps of Antarctica?

  \item What effect does a rough surface, high-resolution (\SI{250}{\metre}) bed topography have on the friction parameters of an ice sheet model?

  \item Where does water drain and accumulate underneath the Antarctic Ice Sheet, how much volume is mobilized, and at what timescales do these processes occur?
\end{enumerate}

The products of this thesis will be a revised continent-wide map of Antarctica's bed topography and subglacial hydrology, including a revised inventory of active subglacial lakes.
Building on top of this, are regional-scale derived products of basal friction and a time-series of ice volume displacement analyzed at select drainage basins.
Taken together, these datasets form a foundation enabling accurate modelling of overland ice flow.

\subsection{Outline}

This thesis is comprised of five chapters.

Chapter 1 establishes the context behind this research, the three research questions, and also the outline you are reading now.
It also contains an introduction to the influence of subglacial topography and subglacial hydrology on ice flow.

Chapter 2 is adapted from a journal paper submitted to The Cryosphere \citep{LeongDeepBedMapdeepneural2020}, reformatted to fit in this thesis.
It starts with an introduction to the field of deep learning in the context of geospatial science.
The chapter then provides a detailed look into the construction of a convolutional neural network architecture to generate a super-resolution (\SI{250}{\metre}) bed elevation map of Antarctica from a low resolution (\SI{1000}{\metre}) BEDMAP2 input and other remote sensing observations of the ice surface.

Chapter 3 uses the neural network generated super-resolution Antarctic bed elevation map to perform an ice sheet model basal inversion over Pine Island Glacier.
In this chapter, we examine the inverted properties of basal drag, effective pressure and friction coefficient using two bed elevation models of different spatial resolutions.
The output of this exercise is to analyze the effects of using a rougher bed on ice flow modelling.

Chapter 4 looks into the active subglacial hydrology system of Antarctica using ICESat-2 laser altimetry.
A revised and automated method for building an active subglacial lake inventory is described, with details on the timing of subglacial lake drainage and filling events, and the estimated volume of water exchanged.
The increased spatiotemporal resolution of these subglacial hydrological maps, combined with an improved subglacial topography, is used to examine the drivers of Antarctic ice flow.

Chapter 5 provides a discussion of the 3 research questions, and implications for future work.
It also presents the main conclusions of this thesis.


\section{Background}

Glaciers and ice sheets move in three main ways:
1) internal deformation of the ice due to gravity
2) sliding of the ice over the basal substrate
3) deformation of the basal substrate.

The latter two processes can be collectively termed as basal slip \citep{Cuffeyphysicsglaciers2010}.
While the bed topography of Antarctica can change significantly over geological time \citep[e.g.][]{HochmuthEvolvingPaleobathymetryCircum2020}, the amount of water available to slide a glacier is more dynamic and can act on sub-annual to annual timescales \citep[e.g.][]{SiegfriedEpisodicicevelocity2016}.
Much focus has been placed on the importance of subglacial water on ice dynamics, such as the draining and filling of subglacial lakes \citep[e.g.][]{Smithinventoryactivesubglacial2009,SiegfriedThirteenyearssubglacial2018} and glacial surges due to rainfall events \citep[e.g.][]{IkenUpliftUnteraargletscherBeginning1983}.
Another contributor is the structure of the basal topography that exerts a drag on the ice-rock contact surface \citep[e.g.][]{Kyrke-SmithRelevanceDetailBasal2018}.
The base of the ice sheet is difficult to observe, as they require oversnow seismic or radar surveys \citep[e.g.][]{HolschuhLinkingpostglaciallandscapes2020} or direct drilling access \citep[e.g.][]{TulaczykWISSARDSubglacialLake2014}.
In contrast, the surface of the ice sheet can be more readily examined using satellites and other remote sensing instruments \citep[e.g.][]{HowatReferenceElevationModel2019,MouginotContinentWideInterferometric2019}.
The study of basal slip at local to regional scales, from sliding due to water and bed deformation, varies with specific conditions at the bed, and can be posed as a geographical question:  % thesis statement
1) Where in Antarctica does bed topography not have an important control on ice flow (i.e. water is more important).
2) Where in Antarctica does bed topography have an important control on ice flow (i.e. a better bed is needed).

To answer this question, we first require an accurate geographic model of the bed, a Digital Elevation Model with sufficient fine-scale structures \citep[Chp.~\ref{ch:2},][]{LeongDeepBedMapdeepneural2020}. Current generations of Antarctic bed elevation models such as BEDMAP2 \citep{FretwellBedmap2improvedice2013} are overly smooth by necessity.
BEDMAP2 is derived from Radio-echo sounding (\gls{RES}) observations, the data was interpolated to a 5 km grid, and then to a 1 km grid.
BedMachine Antarctica \citep{MorlighemMEaSUREsBedMachineAntarctica2020} is generated by mass conservation and is also smoothed in such a way that fails to capture the inherent sub-kilometre scale roughness of an ice sheet's topograhy as seen in paleo-ice sheets uncovered in the modern day.

% TODO make into horizontal plot perhaps
\begin{figure}[htbp]
  \centering
  \begin{subfigure}[t]{0.65\textwidth}
    % TODO get DeepBedMap figure with ground-truth instead of groundtruth
    \includegraphics[width=1.0\textwidth]{chp2_fig6_elevation_roughness_transect}
    \label{fig:1.2a}
  \end{subfigure}
  \vfill
  \begin{subfigure}[t]{0.65\textwidth}
    \includegraphics[width=1.0\textwidth]{chp1_fig_essd-12-2765-2020-f06}
    \label{fig:1.2b}
  \end{subfigure}
  \caption[Comparison of bed elevation products along ground-truth transects]{
    Comparing \gls{RES} bed elevation measurements against different bed products along transect lines in Antarctica.
    % TODO ref{chp:}
    Top panel is over Thwaites Glacier, showing super resolution DeepBedMap\_DEM \citep[purple;][]{LeongDeepBedMapdeepneural2020}, Ground-truth Operation IceBridge \gls{RES} points \citep[orange;][]{ShiMultichannelCoherentRadar2010} and BedMachine Antarctica \citep[green;][]{MorlighemDeepglacialtroughs2019}.
    Bottom panel is over Princess Elizabeth Land, showing ICECAP2 \gls{RES} transects (black), BEDMAP2 (blue), BedMachine (red), and ICECAP2 DEM (green) \citep[Fig. 6 from][]{CuiBedtopographyPrincess2020}.
  }
  \label{fig:1.2}
\end{figure}

The smoothness of existing interpolated bed topography products at the sub-kilometre scale is evident when looking along a transect profile.
In Fig.~\ref{fig:1.2}, we show how fine scale roughness exists in ground-truth radio-echo sounding (\gls{RES}) bed picks.
Continent-wide bed elevation products like BEDMAP2 and BedMachine Antarctica fail to capture this roughness.
Newer regional products over the Weddell Sea sector \citep{Jeofrynewbedelevation2018} and Princess Elizabeth Land \citep{CuiBedtopographyPrincess2020} provide more accurate interpolated bed elevations but are still unable to resolve short-wavelength bed roughness.
Thus, we require a method applicable to the entire continent that retains fine-scale (sub-kilometre) roughness critical for modelling ice flow \citep[see e.g.][]{HubbardSpectralroughnessglaciated2000,SiegertMacroscalebedroughness2004,BinghamDiverselandscapesPine2017,FalciniQuantifyingbedroughness2018}

One way to improve on smooth bed topographies is by using deep learning \citep[see][for a review]{GoodfellowDeeplearning2016}.
Specifically, training a neural network on areas with high resolution datasets, and using it to predict what a high resolution topography would look like in areas with little to no direct observations \citep[Chp.~\ref{ch:2},][]{LeongDeepBedMapdeepneural2020}.
Our deep learning method has similarities with inverse methods \citep{GudmundssonInverseMethodsGlaciology2011}, whereby knowing ice surface elevation, velocity, surface accummulation, and other surface observations in detail allows us to infer what the bed topography might be like.
In contrast to inverse methods, we introduced a specific adversarial loss function \citep{GoodfellowGenerativeAdversarialNetworks2014} that penalizes overtly smooth topography, pushing bed predictions towards that of realistic groundtruth surfaces.

\begin{figure}[htbp]
  \includegraphics[width=0.9\textwidth]{chp1_fig_universal_glacier_slip_law.jpg}
  \caption[Universal glacier slip law]{
    A universal glacier slip law spanning the spectrum of form drag and skin drag regimes.
    Figure is from \citet{Minchewuniversalglacierslip2020}.
  }
  \label{fig:1.3}
\end{figure}

With the DeepBedMap dataset, we then experiment how ice flow over a rough surface differs from ice flow over smooth topography (Chp.~\ref{ch:3}).
Two beds are compared using a Full Stokes ice flow model - a full physical treatment of how an ice stream flows, including the three-dimensional driving forces and stresses that govern the behaviour of an ice body \citep{LarourContinentalscalehigh2012}.
% In addition, we also analyze how water may flow differently over a rough versus smooth topography, whether that has any effect on hydropotential gradients, which may influence basal water pressure at the base of ice streams and thus change ice flow.
% This is crucial to get right, as whether water flows as a concentrated stream in a channel or as a distributed network has a major impact on the ice flowing above it.
One outcome of this exercise is to measure the basal traction of an ice stream, separable into form drag and skin drag \citep[Fig.~\ref{fig:1.3},][]{BinghamDiverselandscapesPine2017,Kyrke-SmithRelevanceDetailBasal2018,Minchewuniversalglacierslip2020,SchoofBasalperturbationsice2002}.
Form drag is basal drag due to topography \citep[c.f.][]{WeertmanSlidingGlaciers1957}, a component which typically increases with higher resolution topography.
Skin drag is the frictional force occurring at the ice-rock interface, determined by bed material properties, and significantly influenced by water that acts to decrease friction \citep{IversonExperimentsdynamicssedimentary2015}.
Skin drag may be overestimated in comparison to form drag on glaciers such as Pine Island Glacier \citep{BinghamDiverselandscapesPine2017,Kyrke-SmithRelevanceDetailBasal2018}, and this detail is of particular relevance when modelling vulnerable glaciers in the Amundsen sea sector \citep{Kyrke-SmithRelevanceDetailBasal2018}.
% TODO explain or throw out. Moreover, it beckons the question on whether basal inversion have been correctly applied to other regions in Antarctica.

% Basal traction itself is normally an inferred property of the bed, and is a standard product produced by ice sheet models where given an ice surface elevation with a known velocity, as well as a bed elevation, we can invert for what basal traction is necessary to produce a desired velocity.
% What is hard to separate out, is the true contribution or division of form drag and skin drag into the basal traction parameter.

Regardless of the nature of subglacial topography, water remains important in determining ice flow.
Water beneath an ice sheet is dynamic, and has been shown to affect the speed of mountain glaciers \citep[e.g.][]{IkenUpliftUnteraargletscherBeginning1983}, and that of outlet glaciers in Greenland where surface water has the ability to access the bed via moulins \citep{ZwallySurfaceMeltInducedAcceleration2002}.
While no active moulins are known to be present in the interior regions of West and East Antarctica \citep{DirscherlAutomatedMappingAntarctic2020,StokesWidespreaddistributionsupraglacial2019}, there is still a considerable volume of water trapped in subglacial lakes \citep{WrightfourthinventoryAntarctic2012}.
Many of these subglacial lakes are active and have drained and filled over the satellite era \citep[Fig.~\ref{fig:1.4},][]{Smithinventoryactivesubglacial2009,SiegfriedThirteenyearssubglacial2018}, including some that were associated with a temporary speed-ups of ice flow \citep[e.g.][]{Scambostriggeringsubglaciallake2011,BellLargesubglaciallakes2007,StearnsIncreasedflowspeed2008,WrightSubglacialhydrologicalconnectivity2014}.

\begin{figure}[htbp]
  \centering
  \begin{subfigure}[t]{1.0\textwidth}
    % TODO get higher resolution image from https://doi.org/10.1017/aog.2017.36 somehow
    \includegraphics[width=1.0\textwidth]{chp1_fig_siegfried_fricker_subglacial_lakes.jpg}
    \label{fig:1.4a}
  \end{subfigure}
  % \vfill
  % \begin{subfigure}[t]{0.8\textwidth}
  %   \includegraphics[width=1.0\textwidth]{chp1_fig_west_antarctica_icesat2_lakes}
  %   \label{fig:1.4b}
  % \end{subfigure}
  \caption[Map of radio-echo sounding and satellite-detected subglacial lakes in Antarctica]{
    Subglacial lakes over Antarctica as detected by satellite and airborne sensors.
    Active subglacial lakes as detected by different satellite sensors are represented by different coloured blobs (blue, light blue and purple).
    Radio-echo sounding (RES) detected subglacial lakes are shown as red blobs.
    % Top panel shows those detected by Cryosat-2 radar altimetry.
    Figure is from \citet{SiegfriedThirteenyearssubglacial2018}.
    % TODO \ref{chp:4}
    % Bottom panel shows those detected by ICESat-2 laser altimetry.
    % Red polygons indicate draining lakes, while blue polygons are filling lakes.
    % Period of subglacial lake activity is from 2018-10-14 to 2020-07-16.
    % Plotted using Antarctic Stereographic Projection (EPSG:3031).
  }
  \label{fig:1.4}
\end{figure}

% Going back to our thesis statement, we need to figure out whether the sensitivity of the Antarctic ice sheet is due to this dynamic water field (skin drag), or if knowing the bed topography (form drag) to a higher resolution itself is sufficient.
% While the study of these two components are not mutually exclusive, and may perhaps be geographically dependent, it is important to be aware of both of them.
% One cannot assume that water in itself is a control on ice flow via sliding (though with subglacial water of sufficient thickness, this may be the case).
% Neither can one assume that a suitably rough bed topography is enough to hold back an ice sheet without an understanding of how water can change the equations in a short time period.

The subglacial hydrological system of Antarctica is very dynamic, and effort is needed to model and account for how subglacial water weakens subglacial till and induce faster sliding.
The Subglacial Hydrology Model Intercomparison Project \citep[SHMIP;][]{DeFleurianSHMIPsubglacialhydrology2018} subjected a land terminating ice sheet to different subglacial hydrological schemes and forcing cycles (steady state, diurnal and seasonal), looking at how effective pressure (ice overburden pressure minus subglacial water pressure) and discharge volumes were affected.
SHMIP found that complex 2-D models which include the physics of drainage systems were needed to model dynamic short term (e.g. diurnal) cycles, while simpler 0-D or 1-D models may suffice for less dynamic scenarios like those in steady state or happening on annual timescales.
In Chapter \ref{ch:4}, our findings of rapid, sub-annual transmission of subglacial water across a connected subglacial lake system suggests that 2-D models are therefore needed.
Our understanding of subglacial hydrology is still inadequate as mentioned in recent model intercomparisons \citep[ISMIP6; ][]{SeroussiISMIP6Antarcticamultimodel2020}, and research will need to focus on coupling subglacial hydrology models with ice sheet models \citep[e.g.][]{SommersSHAKTISubglacialHydrology2018,Smith-Johnsenrolesubglacialhydrology2020} to examine the influence of distributed vs channelized flow on ice dynamics.

One takeaway from this thesis is the need to efficiently integrate data and methods from various fields working on Antarctic glaciology.
In order to inform our understanding of future ice sheet behaviour for sea level rise projections, it is crucial to make full and efficient use of the growing volume of remote sensing and field collected data in our ice sheet models.
The novel machine learning and automated remote sensing methods we introduce present an exciting path forward over existing classical methods, by enabling for the efficient ingestion and merging of diverse datasets.
The bed of Antarctica may remain out of reach for the most part, but one can hope that the increasing synergy between data collectors and modellers will continue to improve our understanding of the Antarctic bed.

% TODO Use motivation at https://www.cambridge.org/core/journals/journal-of-glaciology/article/antarctic-ice-sheet-response-to-sudden-and-sustained-iceshelf-collapse-abumip/C08970BDF95EA737AD941D3690EBB3C5/core-reader#sec5-1-4.
% ``A major uncertainty remains in the physical understanding and modelling of subglacial till mechanics that are essential for determining the effective pressure and sliding rate at the base of a marine ice sheet. The few studies that have attempted to tackle this problem (e.g. Bueler and van Pelt, 2015; Gladstone and others, 2017) are based on seminal work by Tulaczyk and others (2000a, 2000b), but it is clear that more research in subglacial hydrology and basal mechanics of marine ice sheets is needed.''


\chapter{DeepBedMap: a deep neural network for resolving the bed topography of Antarctica}
\label{ch:2}
\chaptermark{DeepBedMap}
% DeepBedMap: a deep neural network for resolving the bed topography of Antarctica

% Macros to follow Copernicus Publications style
% \usepackage{units}
\DeclareRobustCommand*\unit[1]
 {\ensuremath{%
   {\thinmuskip3mu\relax
    \def\mu{\text{\textmu}}\def~{\,}%
    \ifx\f@series\testbx\mathbf{#1}\else\mathrm{#1}\fi}}}
\newcommand\tophline{\hline\noalign{\vspace{1mm}}}
\newcommand\middlehline{\noalign{\vspace{1mm}}\hline\noalign{\vspace{1mm}}}
\newcommand\bottomhline{\noalign{\vspace{1mm}}\hline}

\section*{Abstract}

To resolve the bed elevation of Antarctica, we present DeepBedMap~-- a novel machine learning method that can produce Antarctic bed topography with adequate surface roughness from multiple remote sensing data inputs.
The super-resolution deep convolutional neural network model is trained on scattered regions in Antarctica where high-resolution (250\,\unit{m}) ground-truth bed elevation grids are available.
This model is then used to generate high-resolution bed topography in less surveyed areas.
DeepBedMap improves on previous interpolation methods by not restricting itself to a low-spatial-resolution (1000\,\unit{m}) BEDMAP2 raster image as its prior image.
It takes in additional high-spatial-resolution datasets, such as ice surface elevation, velocity and snow accumulation, to better inform the bed topography even in the absence of ice thickness data from direct ice-penetrating-radar surveys.
The DeepBedMap model is based on an adapted architecture of the Enhanced Super-Resolution Generative Adversarial Network, chosen to minimize per-pixel elevation errors while producing realistic topography.
The final product is a four-times-upsampled (250\,\unit{m}) bed elevation model of Antarctica that can be used by glaciologists interested in the subglacial terrain and by ice sheet modellers wanting to run catchment- or continent-scale ice sheet model simulations.
We show that DeepBedMap offers a rougher topographic profile compared to the standard bicubically interpolated BEDMAP2 and BedMachine Antarctica and envision it being used where a high-resolution bed elevation model is required.

\section{Introduction} %% \introduction[modified heading if necessary]

The bed of the Antarctic ice sheet is one of the most challenging surfaces on Earth to map due to the thick layer of ice cover.
Knowledge of bed elevation is however essential for estimating the volume of ice currently stored in the ice sheets and for input to the numerical models that are used to estimate the contribution ice sheets are likely to make to sea level in the coming century.
The Antarctic ice sheet is estimated to hold a sea level equivalent (SLE) of 57.9\,$\pm$\,0.9\,\unit{m} \citep{MorlighemDeepglacialtroughs2019}.
Between 2012 and 2017, the Antarctic ice sheet was losing mass at an average rate of 219\,$\pm$\,\SI{43}{\giga\tonne\per\year} (0.61\,$\pm$\,\SI{0.12}{\milli\metre\per\year} SLE), with most of the ice loss attributed to the acceleration, retreat and rapid thinning of major West Antarctic Ice Sheet outlet glaciers \citep{IMBIEMassbalanceAntarctic2018}.
Bed elevation exerts additional controls on ice flow by routing subglacial water and providing frictional resistance to flow \citep{SiegertMacroscalebedroughness2004}.
Bed roughness, especially at short wavelengths, exerts a frictional force against the flow of ice, making it an important influence on ice velocity \citep{BinghamDiverselandscapesPine2017,FalciniQuantifyingbedroughness2018}.
The importance of bed elevation has led to major efforts to compile bed elevation models of Antarctica, notably with the BEDMAP1 \citep{LytheBEDMAPnewice2001} and BEDMAP2 \citep{FretwellBedmap2improvedice2013} products.
A need for a higher-spatial-resolution digital elevation model (DEM) is also apparent, as ice sheet models move to using sub-kilometre grids in order to quantify glacier ice flow dynamics more accurately \citep{LeBrocqimprovedAntarcticdataset2010,Grahamhighresolutionsyntheticbed2017}.
Finer grids are especially important at the ice sheet's grounding zone on which adaptive mesh refinement schemes have focused \citep[e.g.][]{CornfordAdaptivemeshrefinement2016}, and attention to the bed roughness component is imperative for proper modelling of fast-flowing outlet glaciers \citep{DurandImpactbedrockdescription2011,NiasContrastingmodelledsensitivity2016}.
Here we address the challenge of producing a high-resolution DEM while preserving a realistic representation of the bed terrain's roughness.

Estimating bed elevation directly from geophysical observations primarily uses ice-penetrating-radar methods \citep[e.g.][]{RobinRadioechoexploration1970}.
Airborne radar methods enable reliable along-track estimates with low uncertainty (around the 1\,{\%} level) introduced by imperfect knowledge of the firn and ice velocity structure, with some potential uncertainty introduced by picking the bed return.
Radar-derived bed estimates remain limited in their geographic coverage \citep{FretwellBedmap2improvedice2013} and are typically anisotropic in their coverage, with higher spatial sampling in the along-track direction than between tracks.

To overcome these limitations, indirect methods of estimating bed elevation have been developed, and these include inverse methods and spatial statistical methods.
Inverse methods use surface observations combined with glaciological-process knowledge to determine ice thickness \citep[e.g.][]{vanPeltiterativeinversemethod2013}.
A non-linear relationship exists between the thickness of glaciers, ice streams and ice sheets and how they flow \citep{Raymondrelationshipsurfacebasal2005}, meaning one can theoretically use a well-resolved surface to infer bed properties \citep[e.g.][]{Farinottimethodestimateice2009}.
Using surface observation inputs, such as the glacier outline, surface digital elevation models, surface mass balance, surface rate of elevation change, and surface ice flow velocity, various models have been tested in the Ice Thickness Models Intercomparison eXperiment \citep[ITMIX;][]{FarinottiHowaccurateare2017} to determine ice thickness (surface elevation minus bed elevation).
While significant inter-model uncertainties do exist, they can be mitigated by combining several models in an ensemble to provide a better consensus estimate \citep{Farinotticonsensusestimateice2019}.
On a larger scale, the inverse technique has also been applied to the Greenland \citep{MorlighemBedMachinev3Complete2017} and Antarctic \citep{MorlighemDeepglacialtroughs2019} ice sheets, specifically using the mass conservation approach \citep{Morlighemmassconservationapproach2011}.
Spatial statistical methods seek to derive a higher-spatial-resolution bed by applying the topographical likeness of bed features known to great detail in one area to other regions.
For example, the conditional simulation method applied by \citet{GoffConditionalsimulationThwaites2014} is able to resolve both fine-scale roughness and channelized morphology over the complex topography of Thwaites Glacier and make use of the fact that roughness statistics are different between highland and lowland areas.
\citet{Grahamhighresolutionsyntheticbed2017} uses a two-step approach to generate their synthetic high-resolution grid, with the high-frequency roughness component coming from the ICECAP and BEDMAP1 compilation radar point data and the low-frequency component coming from BEDMAP2.
Neither method is perfect, and we see all of the above methods as complementary.

We present a deep-neural-network method that is trained on direct ice-penetrating-radar observations over Antarctica and one which has features from both the indirect inverse modelling and spatial statistical methodologies.
An artificial neural network, loosely based on biological neural networks, is a system made up of neurons.
Each neuron comprises a simple mathematical function that takes an input to produce an output value, and neural networks work by combining many of these neurons together.
The term deep neural network is used when there is not a direct function mapping between the input data and final output but two or more layers that are connected to one another \citep[see][for a review]{LeCunDeeplearning2015}.
They are trained using backpropagation, a procedure whereby the weights or parameters of the neurons' connections are adjusted so as to minimize the error between the ground truth and output of the neural network \citep{RumelhartLearningrepresentationsbackpropagating1986}.
Similar work has been done before using artificial neural networks for estimating bed topography\citep[e.g.][]{ClarkeNeuralNetworksApplied2009,MonnierInferencebedtopography2018}, but to our knowledge, no-one so far in the glaciological community has attempted to use convolutional neural networks that work in a more spatially aware, 2-dimensional setting.
Convolutional neural networks differ from standard artificial neural networks in that they use kernels or filters in place of regular neurons \citep[again, see][for a review]{LeCunDeeplearning2015}.
The techniques we employ are prevalent in the computer vision community, having existed since the 1980s \citep{FukushimaNeocognitronnewalgorithm1982,LeCunBackpropagationAppliedHandwritten1989} and are commonly used in visual pattern recognition tasks \citep[e.g.][]{LecunGradientbasedlearningapplied1998,KrizhevskyImageNetClassificationDeep2012}.
Our main contributions are twofold:
we (1)~present a high-resolution (250\,\unit{m}) bed elevation map of Antarctica that goes beyond the 1\,\unit{km} resolution of BEDMAP2 \citep{FretwellBedmap2improvedice2013} and
(2)~design a deep convolutional neural network to integrate as many remote sensing datasets as possible which are relevant to estimating Antarctica's bed topography.
We name the neural network ``DeepBedMap'', and the resulting digital elevation model (DEM) product ``DeepBedMap\_DEM''.


\section{Related Work}

\subsection{Super-Resolution} \label{section:superresolution}

Super resolution involves the processing of a low-resolution raster image into a higher-resolution one \citep{TsaiMultiframeimagerestoration1984}.
The idea is similar to the work on enhancing regular photographs to look crisper.
The problem is especially ill-posed because a specific low-resolution input can correspond to many possible high-resolution outputs, resulting in the development of several different algorithms aimed at solving this challenge \citep[see][for a review]{NasrollahiSuperresolutioncomprehensivesurvey2014}.
One promising approach is to use deep neural networks \citep{LeCunDeeplearning2015} to learn an end-to-end mapping between the low- and high-resolution images, a method coined the Super-Resolution Convolutional Neural Network \citep[SRCNN;][]{DongImageSuperResolutionUsing2014}.
Since the development of SRCNN, multiple advances have been made to improve the perceptual quality of super-resolution neural networks \citep[see][for a review]{YangDeepLearningSingle2019}.
One way is to use a better loss function, also known as a cost function.
A loss function is a mathematical function that represents the error between the output of the neural network and the ground truth (see also Appendix~\ref{appendix:A}).
By having an adversarial component in its loss function, the Super-Resolution Generative Adversarial Network \citep[SRGAN;][]{LedigPhotoRealisticSingleImage2017} manages to produce super-resolution images with finer perceptual details.
A generative adversarial network \citep{GoodfellowGenerativeAdversarialNetworks2014} consists of two neural networks, a generator and a discriminator.
A common analogy used is to treat the generator as an artist that produces imitation paintings and the discriminator as an art critic that determines the authenticity of the paintings.
The artist wants to fool the critic into believing its paintings are real, while the critic tries to identify problems with the painting.
Over time, the artist or generator model learns to improve itself based on the critic's judgement, producing authentic-looking paintings with high perceptual quality.
Perceptual quality is the extent to which an image looks like a valid natural image, usually as judged by a human.
In this case, perceptual quality is quantified mathematically by the discriminator or critic taking into account high-level features of an image like contrast, texture, etc.
Another way to improve performance is by reconfiguring the neural network's architecture, wherein the layout or building blocks of the neural network are changed.
By removing unnecessary model components and adding residual connections \citep{HeDeepResidualLearning2015}, an enhanced deep super-resolution network \citep[EDSR;][]{LimEnhancedDeepResidual2017} features a deeper neural network model that has better performance than older models.
For the DeepBedMap model, we choose to adapt the Enhanced Super-Resolution Generative Adversarial Network \citep[ESRGAN;][]{WangESRGANEnhancedSuperResolution2019} which brings together the ideas mentioned above.
This approach produces state-of-the-art perceptual quality and won the 2018 Perceptual Image Restoration and Manipulation Challenge on Super Resolution (Third Region; \citealp{Blau2018PIRMChallenge2018}).

\subsection{Network Conditioning} \label{section:networkconditioning}

Network conditioning means having a neural network process one source of information in the context of other sources \citep{DumoulinFeaturewisetransformations2018}.
In a geographic context, conditioning is akin to using not just one layer but also other relevant layers with meaningful links to provide additional information for the task at hand.
Many ways exist to insert extra conditional information into a neural network, such as concatenation-based conditioning, conditional biasing, conditional scaling and conditional affine transformations \citep{DumoulinFeaturewisetransformations2018}.
We choose to use the concatenation-based conditioning approach, whereby all of the individual raster images are concatenated together channel-wise, much like the individual bands of a multispectral satellite image.
This was deemed the most appropriate conditioning method as all the contextual remote sensing datasets are raster grid images and also because this approach aligns with related work in the remote sensing field.

An example similar to this DEM super-resolution problem is the classic problem of pan-sharpening, whereby a blurry low-resolution multispectral image conditioned with a high-resolution panchromatic image can be turned into a high-resolution multispectral image.
There is ongoing research into the use of deep convolutional neural networks for pan-sharpening \citep{MasiPansharpeningConvolutionalNeural2016,ScarpaTargetAdaptiveCNNBasedPansharpening2018}, sometimes with the incorporation of specific domain knowledge \citep{YangPanNetDeepNetwork2017}, all of which show promising improvements over classical image processing methods.
More recently, generative adversarial networks \citep{GoodfellowGenerativeAdversarialNetworks2014} have been used in the conditional sense for general image-to-image translation tasks \citep[e.g.][]{IsolaImagetoImageTranslationConditional2016,ParkSemanticImageSynthesis2019}, and also for producing more realistic pan-sharpened satellite images \citep{LiuPSGANGenerativeAdversarial2018}.
Our DeepBedMap model builds upon these ideas and other related DEM super-resolution work \citep{XuNonlocalsimilaritybased2015,ChenConvolutionalNeuralNetwork2016}, while incorporating extra conditional information specific to the cryospheric domain for resolving the bed elevation of Antarctica.


\section{Data and Methods}

\subsection{Data Preparation} \label{section:datapreparation}

Our convolutional neural network model works on 2-D images, so we ensure all the datasets are in a suitable raster grid format.
Ground-truth bed elevation points picked from radar surveys (see Table~\ref{table:groundtruthdata}) are first compiled together onto a common Antarctic stereographic projection (EPSG:3031) using the WGS84 datum, reprojecting where necessary.
These points are then gridded onto a 250\,\unit{m} spatial resolution (pixel-node-registered) grid.
We preprocess the points first using Generic Mapping Tools v6.0 \citep[GMT6;][]{WesselGenericMappingTools2019}, computing the median elevation for each pixel block in a regular grid.
The preprocessed points are then run through an adjustable-tension continuous-curvature spline function with a tension factor set to 0.35 to produce a digital elevation model grid.
This grid is further post-processed to mask out pixels that are more than 3~\unit{pixels} (750\,\unit{m}) from the nearest ground-truth point.

%T1
\begin{table}[t]
  \caption[High-resolution ground-truth datasets for training DeepBedMap]{
    High-resolution ground-truth datasets from ice-penetrating-radar surveys (collectively labelled as~$y$) used to train the DeepBedMap model.
    Training site locations can be seen in Fig.~\ref{fig:2}.
  }
  \label{table:groundtruthdata}
  \begin{tabular}{ll}
  \tophline
  Location                        & Citation                                         \\
  \middlehline
  Pine Island Glacier             & \citet{BinghamDiverselandscapesPine2017}         \\
  Wilkes Subglacial Basin         & \citet{JordanHypothesismegaoutburstflooding2010} \\
  Carlson Inlet                   & \citet{KingIcestreamnot2011}                     \\
  Rutford Ice Stream              & \citet{KingSubglaciallandformsRutford2016}       \\
  Various locations in Antarctica & \citet{ShiMultichannelCoherentRadar2010}         \\
  \bottomhline
  \end{tabular}
\end{table}

%T2
\begin{table*}[t]
  \caption{Remote sensing dataset inputs into the DeepBedMap neural network model.}
  \label{table:datainputs}
  \begin{tabular}{lllll}
  \tophline
  Symbol & Name                        & Variable                                           & Spatial resolution           & Citation                                         \\
  \middlehline
  $x$    & BEDMAP2                     & bed elevation (\unit{m})                           & {~~~~~~~}1000\,\unit{m}               & \citet{FretwellBedmap2improvedice2013}           \\
  $w^1$  & REMA                        & surface elevation (\unit{m})                       & {~~~~~~~~~}100\,\unit{m}$^{\mathrm{b}}$ & \citet{HowatReferenceElevationModel2018}         \\
  $w^2$  & MEaSUREs Ice Velocity       & VX, VY (\unit{m\,yr^{-1}})$^{\mathrm{a}}$          & {~~~~~~~~~}500\,\unit{m}$^{\mathrm{c}}$ & \citet{MouginotContinentwideinterferometric2019} \\
  $w^3$  & Antarctic snow accumulation & snow accumulation rate (\unit{kg\,m^{-2}\,yr^{-1}}) & {~~~~~~~}1000\,\unit{m}               & \citet{ArthernAntarcticsnowaccumulation2006}     \\
  \bottomhline
  \end{tabular}
  \caption*{$^{\mathrm{a}}$\,Note that the $x$ and $y$~components of velocity are used here instead of the norm. \\
  $^{\mathrm{b}}$\,Gaps in 100\,\unit{m} mosaic filled in with bilinear resampled 200\,\unit{m} resolution REMA image. \\
  $^{\mathrm{c}}$\,Originally 450\,\unit{m}; bilinear resampled to 500\,\unit{m}.} % Table Footnotes
\end{table*}

%F1
\begin{figure*}[t]
  \includegraphics[width=170mm]{figures/chp2_fig1_deepbedmap_architecture_compressed}
  \caption[DeepBedMap generator model architecture]{
    DeepBedMap generator model architecture composed of three modules.
    The input module processes each of the four inputs (BEDMAP2, \citealp{FretwellBedmap2improvedice2013}; REMA, \citealp{HowatReferenceElevationModel2019}; MEaSUREs Ice Velocity, \citealp{MouginotMEaSUREsPhaseMap2019}; snow accumulation, \citealp{ArthernAntarcticsnowaccumulation2006}; see also Table~\ref{table:datainputs}) into a consistent tensor.
    The core module processes the rich information contained within the concatenated inputs.
    The upsampling module scales the tensor up by 4 times and does some extra processing to produce the output DeepBedMap\_DEM.
  }
  \label{fig:1}
\end{figure*}

To create the training dataset, we use a sliding window to obtain square tiles cropped from the high-resolution (250\,\unit{m}) ground-truth bed elevation grids, with each tile required to be completely filled with data (i.e.
no Not a Number -- NaN -- values).
Besides these ground-truth bed elevation tiles, we also obtain other tiled inputs (see Table~\ref{table:datainputs}) corresponding to the same spatial bounding box area.
To reduce border edge artefacts in the prediction, the neural network model's input convolutional layers (see Fig.~\ref{fig:1}) use no padding (also known as ``valid'' padding) when performing the initial convolution operation.
This means that the model input grids ($x$, $w^1$, $w^2$, $w^3$) have to cover a larger spatial area than the ground-truth grids~($y$).
More specifically, the model inputs cover an area of 11\,\unit{km}\,$\times$\,11\,\unit{km} (e.g. 11~\unit{pixels}\,$\times$\,11~\unit{pixels} for BEDMAP2), while the ground-truth grids cover an area of 9\,\unit{km}\,$\times$\,9\,\unit{km} (36~\unit{pixels}\,$\times$\,36~\unit{pixels}).
As the pixels of the ground-truth grids may not align perfectly with those of the model's input grids, we use bilinear interpolation to ensure that all the input grids cover the same spatial bounds as those of the reference ground-truth tiles.
The general locations of these training tiles are shown in orange in Fig.~\ref{fig:2}.

\subsection{Model Design} \label{section:modeldesign}

Our DeepBedMap model is a generative adversarial network \citep{GoodfellowGenerativeAdversarialNetworks2014} composed of two convolutional neural network models, a generator $G_\theta$ that produces the bed elevation prediction and a discriminator $D_\eta$ critic that will judge the quality of this output.
The two models are trained to compete against each other, with the generator trying to produce images that are misclassified as real by the discriminator and the discriminator learning to spot problems with the generator's prediction in relation to the ground truth.
Following this is a mathematical definition of the neural network models and their architecture.

The objective of the main super-resolution generator model $G_\theta$ is to produce a high-resolution (250\,\unit{m}) grid of Antarctica's bed elevation $\hat{y}$ given a low-resolution (1000\,\unit{m}) BEDMAP2 \citep{FretwellBedmap2improvedice2013} image~$x$.
However, the information contained in BEDMAP2 is insufficient for this regular super-resolution task, so we provide the neural network with more context through network conditioning (see Sect.~\ref{section:networkconditioning}).
Specifically, the model is conditioned at the input block stage with three raster grids (see Table~\ref{table:datainputs}): (1)~ice surface elevation~$w^1$, (2)~ice surface velocity~$w^2$ and (3)~snow accumulation~$w^3$.
This can be formulated as follows:

\begin{equation}\label{eq:1}
  \hat{y} = G_\theta(x, w^1, w^2, w^3),
\end{equation}

where $G_\theta$ is the generator (see Fig.~\ref{fig:1}) that produces high-resolution image candidates~$\hat{y}$.
For brevity in the following equations, we simplify Eq.~(\ref{eq:1}) to hide conditional inputs $w^1, w^2$ and  $w^3$, so that all input images are represented using~$x$.
To train the generative adversarial network, we update the parameters of the generator~$\theta$ and discriminator~$\eta$ as follows:

\begin{align}
  &\hat{\theta} = \arg\min_{\theta} \frac{1}{N}\sum_{n=1}^{N}L_{\mathrm{G}}(\hat{y}_n, y_n), \label{eq:2}\\
  &\hat{\eta} = \arg\min_{\eta} \frac{1}{N}\sum_{n=1}^{N}L_{\mathrm{D}}(\hat{y}_n, y_n), \label{eq:3}
\end{align}

where new estimates of the neural network parameters~$\hat{\theta}$ and~$\hat{\eta}$ are produced by minimizing the total loss functions $L_{\mathrm{G}}$ and $L_{\mathrm{D}}$, respectively, for the generator~$G$ and discriminator~$D$ and $\hat{y}_n$ and $y_n$ are the set of predicted and ground-truth high-resolution images over $N$~training samples.
The generator network's loss~$L_{\mathrm{G}}$ is a custom perceptual loss function with four weighted components~-- content, adversarial, topographic and structural loss.
The discriminator network's loss~$L_{\mathrm{D}}$ is designed to maximize the likelihood that predicted images are classified as fake~(0) and ground-truth images are classified as real~(1).
Details of these loss functions are described in Appendix~\ref{appendix:A}.

Noting that the objective of the Generator $G$ is opposite to that of the Discriminator $D$, we formulate the adversarial min-max problem following \citet{GoodfellowGenerativeAdversarialNetworks2014} as so:

\begin{equation}\label{eq:4}
  \begin{split}
  \min_{G}\,\max_{D} V(G,D) =&~\mathbb{E}_{y \sim P_{\text{data}}(y)}[\ln D(y)]\\
  &+ \mathbb{E}_{x \sim P_{G(x)}}[\ln(1-D(G(x)))],
  \end{split}
\end{equation}

where for the discriminator~$D$, we maximize the expectation~$\mathbb{E}$ or the likelihood that the probability distribution of the discriminator's output fits $D(y)=1$ when $y \sim P_{\text{data}}(y)$; i.e. we want the discriminator to classify the high-resolution image as real (1) when the image~$y$ is in the distribution of the ground-truth images $P_{\text{data}}(y)$.
For the generator~$G$, we minimize the likelihood that the discriminator classifies the generator output $D(G(x))=0$ when $x \sim P_{G(x)}$; i.e. we do not want the discriminator to classify the super-resolution image as fake~(0) when the inputs~$x$ are in the distribution of generated images~$P_{G(x)}$.
The overall goal of the entire network is to make the distribution of generated images~$G(x)$ as similar as possible to the ground truth~$y$ through optimizing the value function~$V$.

%F2
\begin{figure*}[t]
    \includegraphics[width=120mm]{figures/chp2_fig2_deepbedmap_dem_compressed}
    \caption[DeepBedMap\_DEM overview map over the entire Antarctic continent]{
      DeepBedMap\_DEM over the entire Antarctic continent.
      Plotted on an Antarctic stereographic projection (EPSG:3031) with elevation referenced to the WGS84 datum.
      Grounding line is plotted as thin black line.
      Purple box shows Pine Island Glacier extent used in Fig.~\ref{fig:3}.
      Yellow box shows Thwaites Glacier extent used in Fig.~\ref{fig:5}.
      Orange areas show locations of training tiles (see Table~\ref{table:groundtruthdata}).
    }
    \label{fig:2}
\end{figure*}

DeepBedMap's model architecture is adapted from the Enhanced Super-Resolution Generative Adversarial Network
\citep[ESRGAN;][]{WangESRGANEnhancedSuperResolution2019}.
The generator model~$G$ (see Fig.~\ref{fig:1}) consists of an input, core and upsampling module.
The input module is made up of four sub-networks, each one composed of a convolutional neural network that processes the input image into a consistent 9\,$\times$\,9 shaped tensor.
Note that the MEaSUREs Ice Velocity \citep{MouginotMEaSUREsPhaseMap2019} input has two channels, one each for the $x$ and $y$~velocity components.
All the processed inputs are then concatenated together channel-wise before being fed into the core module.
The core module is based on the ESRGAN architecture with 12 residual-in-residual dense blocks \citep[see][for details]{WangESRGANEnhancedSuperResolution2019}, saddled in between a pre-residual and post-residual convolutional layer.
A skip connection runs from the pre-residual layer's output to the post-residual layer's output before being fed into the upsampling module.
This skip connection \citep{HeIdentityMappingsDeep2016} helps with the neural network training process by allowing the model to also consider minimally processed information from the input module, instead of solely relying on derived information from the residual-block layers when performing the upsampling.
The upsampling module is composed of two upsampling blocks, specifically a nearest-neighbour upsampling followed by a convolutional layer and leaky rectified linear unit \citep[LeakyReLU;][]{MaasRectifiernonlinearitiesimprove2013} activation, which progressively scales the tensors by 2 times each time.
Following this are two deformable convolutional layers \citep{DaiDeformableConvolutionalNetworks2017} which produce the final-output super-resolution DeepBedMap\_DEM.
This generator model is trained to gradually improve its prediction by comparing the predicted output with ground-truth images in the training regions (see Fig.~\ref{fig:2}), using the total loss function defined in Eq.~(\ref{eq:A9}).

The main differences between the DeepBedMap generator model and ESRGAN are the custom input block at the beginning and the deformable convolutional layers at the end.
The custom input block is designed to handle the prior low-resolution BEDMAP2 image and conditional inputs (see Table~\ref{table:datainputs}).
Deformable convolution was chosen in place of the standard convolution so as to enhance the model's predictive capability by having it learn dense spatial transformations.

Besides the generator model, there is a separate adversarial discriminator model $D$ (not shown in the paper).
Again, we follow ESRGAN's \citep{WangESRGANEnhancedSuperResolution2019} lead by implementing the adversarial discriminator network in the style of the Visual Geometry Group convolutional neural network model \citep[VGG;][]{SimonyanVeryDeepConvolutional2014}.
The discriminator model consists of 10 blocks made up of a convolutional, batch normalization \citep{IoffeBatchNormalizationAccelerating2015} and LeakyReLU \citep{MaasRectifiernonlinearitiesimprove2013} layer, followed by two fully connected layers comprised of 100\,\unit{neurons} and 1\,\unit{neuron}, respectively.
For numerical stability, we omit the final fully connected layer's sigmoid activation function from the discriminator model's construction, integrating it instead into the binary cross-entropy loss functions at Eqs.~(\ref{eq:A2}) and~(\ref{eq:A3}) using the log-sum-exp function.
The output of this discriminator model is a value ranging from~0 (fake) to~1 (real) that scores the generator model's output image.
This score is used by both the discriminator and generator in the training process and helps to push the predictions towards more realistic bed elevations.
More details of the neural network training setup can be found in Appendix~\ref{appendix:B}.


\section{Results}

\subsection{DeepBedMap\_DEM Topography} \label{section:deepbedmapdemtopography}

Here we present the output digital elevation model (DEM) of the super-resolution DeepBedMap neural network model and compare it with bed topography produced by other methods.
The resulting DEM has a 250\,\unit{m} spatial resolution and therefore a four-times upsampled bed elevation grid product of BEDMAP2 \citep{FretwellBedmap2improvedice2013}.
In Fig.~\ref{fig:2}, we show that the full Antarctic-wide DeepBedMap\_DEM manages to capture general topographical features across the whole continent.
The model is only valid for grounded-ice regions, but we have produced predictions extending outside of the grounding-zone area (including ice shelf cavities) using the same bed elevation, surface elevation, ice velocity and snow accumulation inputs where such data are available up to the ice shelf front.
We emphasize that the bed elevation under the ice shelves has not been super resolved properly and is not intended for ice sheet modelling use.
Users are encouraged to cut the DeepBedMap\_DEM using their preferred grounding line \citep[e.g.][]{BindschadlerGettingAntarcticanew2011,RignotAntarcticgroundingline2011,MouginotMEaSURESAntarcticBoundaries2017} and replace the under-ice-shelf areas with another bathymetry grid product \citep[e.g.][]{GEBCOCompilationGroupGEBCO2020Grid2020}.
The transition from the DeepBedMap\_DEM to the bathymetry product across the grounding zone can then be smoothed using inverse distance weighting or an alternative interpolation method.

%F3
\begin{figure*}[t]
  \includegraphics[width=130mm]{figures/chp2_fig3_qualitative_bed_comparison}
  \caption[2-D comparison of interpolated bed elevation grid products over Pine Island Glacier]{
    Comparison of interpolated bed elevation grid products over Pine Island Glacier (see extent in Figure \ref{fig:2}).
    \textbf{a} DeepBedMap (ours) at 250 m resolution.
    \textbf{b} BEDMAP2 \citep{FretwellBedmap2improvedice2013}, originally 1000 m, bicubic interpolated to 250 m.
    \textbf{c} Elevation Difference between DeepBedMap and BEDMAP2.
    \textbf{d} BedMachine Antarctica \citep{MorlighemMEaSUREsBedMachineAntarctica2019}, originally 500 m, bicubic interpolated to 250 m.
  }
  \label{fig:3}
\end{figure*}

%F4
\begin{figure*}[t]
  \includegraphics[width=170mm]{figures/chp2_fig4_deepbedmap_closeups}
  \caption[2-D Close-up views of DeepBedMap\_DEM over Antarctica]{
    Close-up views of DeepBedMap\_DEM around Antarctica.
    Panels \textbf{(a--c)}~show Siple Coast locations.
    Panels \textbf{(d--f)}~show Weddell Sea region locations.
    Panels \textbf{(g--i)}~show East Antarctica locations.
    Features of interest are annotated in black text against a white background:
    ridges~R, speckle patterns~S, terraces~T, wave patterns~W.
  }
  \label{fig:4}
\end{figure*}

We now highlight some qualitative observations of DeepBedMap\_DEM's bed topography beneath Pine Island Glacier (Figure \ref{fig:3}) and other parts of Antarctica (Figure \ref{fig:4}).
DeepBedMap\_DEM shows a terrain with realistic topographical features, having fine-scale bumps and troughs that makes it rougher than that of BEDMAP2 \citep{FretwellBedmap2improvedice2013} and BedMachine Antarctica \citep{MorlighemMEaSUREsBedMachineAntarctica2019} while still preserving the general topography of the area (Figure \ref{fig:3}).
Over steep topographical areas such as the Transantarctic Mountains (Figure \ref{fig:4}a, \ref{fig:4}h), DeepBedMap produced speckle (\textbf{S}) texture patterns.
Along fast flowing ice streams and glaciers (Figure \ref{fig:4}b, \ref{fig:4}c, \ref{fig:4}d, \ref{fig:4}e, \ref{fig:4}f, \ref{fig:4}g, \ref{fig:4}h), we can see ridges (\textbf{R}) aligned parallel to the sides of the valley, i.e. along flow.
In some cases, the ridges are also oriented perpendicular to the flow direction such at Whillans Ice Stream (Figure \ref{fig:4}b), Bindschadler Ice Stream (Figure \ref{fig:4}c) and Totten Glacier (Figure \ref{fig:4}g), resulting in intersecting ridges that creates a box-like, honeycomb structure.
Over relatively flat regions in both West and East Antarctica (e.g. Figure \ref{fig:4}g), there are some hummocky, wave-like (\textbf{W}) patterns occasionally represented in the terrain.
Terrace (\textbf{T}) features can occasionally be found winding along the side of hills such as at the Gamburtsev Subglacial Mountains (Figure \ref{fig:4}i).


\subsection{Surface Roughness} \label{section:surfaceroughness}

%F5
\begin{figure*}[t]
  \includegraphics[width=130mm]{figures/chp2_fig5_elevation_roughness_grids}
  \caption[Spatial 2-D view of bed elevation and roughness grids over Thwaites Glacier]{
    Spatial 2-D view of grids over Thwaites Glacier, West Antarctica.
    Plotted on an Antarctic stereographic projection (EPSG:3031) with elevation and SD values in metres referenced to the WGS84 datum.
    \textbf{(a)}~DeepBedMap digital elevation model.
    \textbf{(b)}~2-D roughness from the DeepBedMap\_DEM grid.
    \textbf{(c)}~2-D roughness from interpolated Operation IceBridge grid.
    \textbf{(d)}~2-D roughness from bicubically interpolated BedMachine Antarctica grid.
    Orange points in~\textbf{(a)} correspond to transect sampling locations used in Fig.~\ref{fig:6}.
  }
  \label{fig:5}
\end{figure*}

We compare the roughness of DeepBedMap\_DEM vs. BedMachine Antarctica with ground-truth grids from processed Operation IceBridge data \citep{ShiMultichannelCoherentRadar2010} using SD as a simple measure of roughness \citep{RippinBasalroughnessInstitute2014}.
We calculate the surface roughness for a single 250\,\unit{m} pixel from the SD of elevation values over a square 1250\,\unit{m}\,$\times$\,1250\,\unit{m} area (i.e. 5~\unit{pixels}\,$\times$\,5~\unit{pixels}) surrounding the central pixel.
Focusing on Thwaites Glacier, the spatial 2-D view of the DeepBedMap\_DEM (Fig.~\ref{fig:5}a) shows a range of typical topographic features such as hills and canyons.
The calculated 2-D roughnesses for both DeepBedMap\_DEM (Fig.~\ref{fig:5}b) and the Ground truth (Fig.~\ref{fig:5}c) lie in a similar range from 0 to 400\,\unit{m}, whereas the roughness of BedMachine Antarctica (Fig.~\ref{fig:5}d) is mostly in the 0-to-200\,\unit{m} range (hence the different colour scale).
Also, the roughness pattern for both DeepBedMap\_DEM and the ground truth has a more distributed cluster pattern made up of little pockets (especially towards the coastal region on the left; see Fig.~\ref{fig:5}b and~c), whereas the BedMachine Antarctica roughness pattern shows larger cluster pockets in isolated regions (see Fig.~\ref{fig:5}d).

Taking a 1-D transect over the 250\,\unit{m} resolution DeepBedMap\_DEM, BedMachine Antarctica and ground-truth grids, we illustrate the differences in bed topography and roughness from the coast towards the inland area of Thwaites Glacier with a flight trace from Operation IceBridge (see Fig.~\ref{fig:6}).
For better comparison, we have calculated the Operation IceBridge ground-truth bed elevation and roughness values from a resampled 250\,\unit{m} grid instead of using its native along-track resolution.
All three elevation profiles are shown to follow the same general trend from the relatively rough coastal region (Fig.~\ref{fig:6}a from $-$1550 to $-$1500\,\unit{km} on the $x$~scale), along the retrograde slope (Fig.~\ref{fig:6}a from $-$1500 to $-$1450\,\unit{km} on the $x$~scale) and into the interior region.
DeepBedMap\_DEM features a relatively noisy elevation profile with multiple fine-scale ($<$\,10\,\unit{km}) bumps and troughs similar to the ground truth, while BedMachine Antarctica shows a smoother profile that is almost a moving average of the ground-truth elevation (Fig.~\ref{fig:6}a).
Looking at the roughness statistic (Fig.~\ref{fig:6}b), both the DeepBedMap\_DEM and Operation IceBridge ground-truth grids have a mean SD of about 40\,\unit{m}, whereas BedMachine Antarctica has a mean of about 10\,\unit{m} and rarely exceeds a SD value of 20\,\unit{m} along the transect.

%F6
\begin{figure}[t]
  \includegraphics[width=83mm]{figures/chp2_fig6_elevation_roughness_transect}
  \caption[Transect plots of bed elevation and surface roughness of interpolated grid products]{
    Comparing bed elevation~\textbf{(a)} and surface roughness~\textbf{(b)} (SD of elevation values) of each interpolated grid product (250\,\unit{m} resolution) over a transect (see Fig.~\ref{fig:5} for location of transect line).
    Purple values are from the super-resolution DeepBedMap\_DEM;
    orange values are from tension-spline-interpolated Operation IceBridge ground-truth points;
    green values are from bicubically interpolated BedMachine Antarctica.
  }
  \label{fig:6}
\end{figure}


\section{Discussion}

\subsection{Bed Features}

In Sect.~\ref{section:deepbedmapdemtopography}, we show that the DeepBedMap model has produced a high-resolution (250\,\unit{m}) result (see Fig.~\ref{fig:3}) that can capture a detailed picture of the underlying bed topography.
The fine-scale bumps and troughs are the result of the DeepBedMap generator model learning to produce features that are similar to those found in the high-resolution ground-truth datasets it was trained on.
However, there are also artefacts produced by the model.
For example, the winding terrace (T, Fig.~\ref{fig:4}) features are hard to explain, and though they resemble eskers \citep{DrewsActivelyevolvingsubglacial2017}, their placement along the sides of hills does not support this view.
Similarly, we are not sure why speckle (S, Fig.~\ref{fig:4}) texture patterns are found over steep mountains, but the lack of high-resolution training datasets likely leads the model to perform worse over these high-gradient areas.

Another issue is that DeepBedMap will often pick up details from the high-resolution ice surface elevation model \citep{HowatReferenceElevationModel2019} input dataset, which may not be representative of the true bed topography.
For example, the ridges (R, Fig.~\ref{fig:4}) found along fast-flowing ice streams and glaciers are likely to be the imprints of crevasses or flow stripes \citep{GlasserLongitudinalsurfacestructures2012} observable from the surface.
An alternative explanation is that the ridges, especially the honeycomb-shaped ones, are rhombohedral moraine deposits formed by soft sediment squeezed up into basal crevasses that are sometimes found at stagnant surging glaciers \citep{Dowdeswellvarietydistributionsubmarine2016,DowdeswellRhombohedralcrevassefillridges2016,SolheimSeafloormorphologyoutside1985}.
We favour the first interpretation as the positions of these bed features coincide with the surface features and also because these ridges are more likely to be eroded away in these fast-flowing ice stream areas.

% \hack{\newpage}

The hummocky wave-like (W) patterns we observe over the relatively flat and slower-flowing areas are likely to result from surface megadune structures \citep{ScambosSnowMegadune2014}.
Alternatively, they may be ribbed or Rogen moraine features that are formed in an orientation transverse to the ice flow direction \citep{HattestrandRibbedmorainesSweden1997,HattestrandRibbedmoraineformation1999}.
While any one of these two explanations may be valid in different regions of Antarctica, we lean towards the conservative interpretation that these features are the result of the DeepBedMap model overfitting to the ice surface elevation data.

\subsection{Roughness}

In Sect.~\ref{section:surfaceroughness}, we quantitatively show that a well-trained DeepBedMap neural network model can produce high roughness values more comparable to the ground truth than those of BedMachine Antarctica.
While the mass conservation technique used by BedMachine Antarctica \citep{MorlighemDeepglacialtroughs2019} improves upon ordinary interpolation techniques such as bicubic interpolation and kriging, its results are still inherently smooth by nature.
The ground-truth grids show that rough areas do exist on a fine scale, and so the high-resolution models we produce should reflect that.

DeepBedMap\_DEM manages to capture much of the rough topography found in the Operation IceBridge ground-truth data, especially near the coast (see Fig.~\ref{fig:6}a, from $-$1550 to $-$1500\,\unit{km} on the $x$~scale) where the terrain tends to be rougher.
Along the retrograde slope (see Fig.~\ref{fig:6}a, from $-$1500 to $-$1450\,\unit{km} on the $x$~scale), several of the fine-scale ($<$\,10\,\unit{km}) bumps and troughs in DeepBedMap\_DEM can be seen to correlate well in position with the ground truth.
In contrast, the cubically interpolated BedMachine Antarctica product lacks such fine-scale ($<$\,10\,\unit{km}) bumps and troughs, appearing as a relatively smooth terrain over much of the transect.
Previous studies that estimated basal shear stress over Thwaites Glacier have found a band of strong bed extending about 80--100\,\unit{km} from the grounding line, with pockets of weak bed interspersed between bands of strong bed further upstream \citep{JoughinBasalconditionsPine2009,SergienkoRegularPatternsFrictional2013}, a pattern that is broadly consistent with the DeepBedMap\_DEM roughness results (see Fig.~\ref{fig:5}b).

In general, DeepBedMap\_DEM produces a topography that is rougher, with SD values more in line with those observed in the ground truth (see Fig.~\ref{fig:6}b).
The roughness values for BedMachine Antarctica are consistently lower throughout the transect, a consequence of the mass conservation technique using regularization parameters that yield smooth results.
We note that the DeepBedMap\_DEM does appear rougher than the ground truth in certain areas.
It is possible to tweak the training regime to incorporate roughness (or any statistical measure) into the loss function (see Appendix~\ref{appendix:A}) to yield the desired surface, and this will be explored in future work (see Sect.~\ref{section:futuredirections}).
Recent studies have stressed the importance of form drag (basal drag due to bed topography) over skin drag (or basal friction) on the basal traction of Pine Island Glacier \citep{BinghamDiverselandscapesPine2017,Kyrke-SmithRelevanceDetailBasal2018}, and the DeepBedMap super-resolution work here shows strong potential in meeting that demand as a high-resolution bed topography dataset for ice sheet modelling studies.

In terms of bed roughness anisotropy, DeepBedMap is able to capture aspects of it from the ground-truth grids by combining (1)~ice flow direction via the ice velocity grid's $x$ and $y$~components \citep{MouginotMEaSUREsPhaseMap2019}, (2)~ice surface aspect via the ice surface elevation grid \citep{HowatReferenceElevationModel2019}, and (3)~the low-resolution bed elevation input \citep{FretwellBedmap2improvedice2013}.
There are therefore inherent assumptions that the topography of the current bed is associated with the current ice flow direction, surface aspect and existing low-resolution BEDMAP2 anisotropy.
Provided that the direction of this surface velocity and aspect is the same as bed roughness anisotropy, as demonstrated in \citet{HolschuhLinkingpostglaciallandscapes2020}, the neural network will be able to recognize it and perform accordingly.
However, if the ice flow direction and surface aspect is not associated with bed anisotropy, then this assumption will be violated and the model will not perform well.

\subsection{Limitations}

The DeepBedMap model is trained only on a small fraction of the area of Antarctica, at less than 0.1\,{\%} of the grounded-ice regions (excluding ice shelves and islands).
This is because the pixel-based convolutional neural network cannot be trained on sparse survey point measurements, nor is it able to constrain itself with track-based radar data.
As the along-track resolution of radar bed picks are much smaller than 250\,\unit{m} pixels, it is also not easy to preserve roughness from radar unless smaller pixels are used.
The topography generated by the model is sensitive to the accuracy of its data inputs (see Tables~\ref{table:groundtruthdata} and \ref{table:datainputs}), and though this is a problem faced by other inverse methods, neural network models like the one presented can be particularly biased towards the training dataset.
Specifically, the DeepBedMap model focuses on resolving short-wavelength features important for sub-kilometre roughness, compared to BedMachine Antarctica \citep{MorlighemDeepglacialtroughs2019} which recovers large-scale features like ridges and valleys well.

An inherent assumption in this methodology is that the training datasets have sampled the variable bed lithology of Antarctica \citep{CoxGeoMAPdatasetAntarctic2018} sufficiently.
This is unlikely to be true, introducing uncertainty into the result as different lithologies may cause the same macroscale bed landscapes to result in a range of surface features.
In particular, the experimental model's topography is likely skewed towards the distribution of the training regions that tend to reside in coastal regions, especially over ice streams in West Antarctica (see Fig.~\ref{fig:2}).
While bed lithology could be used as an input to inform the DeepBedMap model's prediction, it is challenging to find a suitable geological map (or geopotential proxy; \citealp[see e.g.][]{AitkensubglacialgeologyWilkes2014,CoxGeoMAPdatasetAntarctic2018}) for the entire Antarctic continent that has a sufficiently high spatial resolution.
Ideally, the lithological map (categorical or qualitative) would first be converted to a hardness map with an appropriate erosion law and history incorporated (quantitative).
This is because it is easier to train generative adversarial networks on quantitative data (e.g. hardness as a scale from 0 to 10) than on categorical data variables (e.g. sedimentary, igneous or metamorphic rocks); the latter would require a more elaborate model architecture and loss function design.

\subsection{Future directions} \label{section:futuredirections}

The way forward for DeepBedMap is to combine quality datasets gathered by radioglaciology and remote sensing specialists, with new advancements made by the ice sheet modelling and machine learning community.
While care has been taken to source the best possible datasets (see Tables~\ref{table:groundtruthdata} and~\ref{table:datainputs}), we note that there are still areas where more data are needed.
Radio-echo sounding is the best tool available to fill in the data gap, as it provides not only the high-resolution datasets needed for training but also the background coarse-resolution BEDMAP dataset.
Besides targeting radio-echo-sounding acquisitions over a diverse range of bed and flow types, swath reprocessing of old datasets that have that capability \citep{HolschuhLinkingpostglaciallandscapes2020} may be another useful addition to the training set.
The super-resolution DeepBedMap technique can also be applied on bed elevation inputs newer than BEDMAP2 \citep{FretwellBedmap2improvedice2013}, such as the 1000\,\unit{m} resolution DEM over the Weddell Sea \citep{Jeofry1KmBedTopography2017}, the 500\,\unit{m} resolution BedMachine Antarctica dataset \citep{MorlighemMEaSUREsBedMachineAntarctica2019} or the upcoming BEDMAP3.

A way to increase the number of high-resolution ground-truth training data further is to look at formerly glaciated beds.
There are a wealth of data around the margins of Antarctica in the form of swath bathymetry data and also on land in areas like the former Laurentide ice sheet.
The current model architecture does not support using solely ``elevation'' as an input, because it also requires ice elevation, ice surface velocity and snow accumulation data.
In order to support using these paleobeds as training data, one could do one of the following:

\begin{enumerate}
  \item Have a paleo-ice-sheet model that provides these ice surface observation parameters.
  However, continent-scale ice sheet models quite often produce only kilometre-scale outputs, and there are inherent uncertainties with past ice sheet reconstructions that may bias the resulting trained neural network model.

  \item Modularize the neural network model to support different sets of training data.
  One main branch would be trained like a single-image super-resolution problem \citep{YangDeepLearningSingle2019}, where we try to map a low-resolution BEDMAP2 tile to a high-resolution ground-truth image (be it from a contemporary bed, a paleobed or offshore bathymetry).
  The optional conditional branches would then act to support and improve on the result of this naive super-resolution method.
  This design is more complicated to set up and train, but it can increase the available training data by at least an order of magnitude and lead to better results.
\end{enumerate}

From a satellite remote sensing perspective, it is important to continue the work on increasing spatial coverage and measurement precision.
Some of the conditional datasets used such as REMA \citep{HowatReferenceElevationModel2019} and MEaSUREs Ice Velocity \citep{MouginotMEaSUREsPhaseMap2019} contain data gaps which introduce artefacts in the DeepBedMap\_DEM, and those holes need to be patched up for proper continent-wide prediction.
A surface mass balance dataset with sub-kilometre spatial resolution will also prove useful in replacing the snow accumulation dataset \citep{ArthernAntarcticsnowaccumulation2006} used in this work.
As the DeepBedMap model relies on data from multiple sources collected over different epochs, it has no proper sense of time.
Ice elevation change captured using satellite altimeters such as from CryoSat-2 \citep{HelmElevationelevationchange2014}, ICESat-2 \citep{MarkusIceCloudland2017} or the upcoming CRISTAL \citep{KernCopernicusPolarIce2020} could be added as an additional input to better account for temporal factors.

The DeepBedMap model's modular design (see Sect.~\ref{section:modeldesign}) means the different modules (see Fig.~\ref{fig:1}) can be improved on and adapted for future-use cases.
The generator model architecture's input module can be modified to handle new datasets such as the ones suggested above or redesigned to extract a greater amount of information for better performance.
Similarly, the core and upsampling modules which are based on ESRGAN \citep{WangESRGANEnhancedSuperResolution2019} can be replaced with newer, better architectures as the need arises.
The discriminator model which is currently one designed for standard computer vision tasks can also be modified to incorporate glaciology-specific criteria.
For example, the generated bed elevation image could be scrutinized by the discriminator model for valid properties such as topographic features that are aligned with the ice flow direction.
The redesigned neural network model can be retrained from scratch or fine-tuned using the trained weights from DeepBedMap to further improve the predictive performance.
Taken together, these advances will lead to an even more accurate and higher-resolution bed elevation model.


\section{Conclusions}  %% \conclusions[modified heading if necessary]

The DeepBedMap convolutional neural network method presents a data-driven approach to resolve the bed topography of Antarctica using existing data.
It is an improvement beyond simple interpolation techniques, generating high-spatial-resolution (250\,\unit{m}) topography that preserves detail in bed roughness and is adaptable for catchment- to continent-scale studies on ice sheets.
Unlike other inverse methods that rely on some explicit parameterization of ice flow physics, the model uses deep learning to find suitable neural network parameters via an iterative error minimization approach.
This makes the resulting model particularly sensitive to the training dataset, emphasizing the value of densely spaced bed elevation datasets and the need for such sampling over a more diverse range of Antarctic substrate types.
The use of graphical processing units (GPUs) for training and inference allows the neural network method to scale easily, and the addition of more training datasets will allow it to perform better.

The work here is intended not to discourage the usage of other inverse modelling or spatial statistical techniques but to introduce an alternative methodology, with an outlook towards combining each methodology's strengths.
Once properly trained, the DeepBedMap model runs quickly (about 3\,min for the whole Antarctic continent) and produces realistic rough topography.
Combining the DeepBedMap model with more physically based mass conservation inverse approaches \citep[e.g.][]{MorlighemDeepglacialtroughs2019} will likely result in more efficient ways of generating accurate bed elevation maps of Antarctica.
One side product resulting from this work is a test-driven development framework that can be used to measure and compare the performance of upcoming bed terrain models.
The radioglaciology community has already begun to compile a new comprehensive bed elevation and ice thickness dataset for Antarctica, and there have been discussions on combining various terrain interpolation techniques in an ensemble to collaboratively create the new BEDMAP3.



\begin{appendices}

\section{Details of loss function components} \label{appendix:A}

The loss function, or cost function, is a mathematical function that maps a set of input variables to an output loss value.
The loss value can be thought of as a weighted sum of several error metrics between the neural network's prediction and the expected output or ground truth.
It is this loss value which we want to minimize so as to train the neural network model to perform better, and we do this by iteratively optimizing the parameters in the loss function.
Following this are the details of the various loss functions that make up the total loss function of the DeepBedMap generative adversarial network model.

\subsection{Content Loss}

To bring the pixel values of the generated images closer to those of the ground truth, we first define the content-loss function~$L_1$.
Following ESRGAN \citep{WangESRGANEnhancedSuperResolution2019}, we have

\begin{equation}\label{eq:A1}
  L_1 = \dfrac{1}{n} \sum\limits_{i=1}^n ||\hat{y}_i - y_i||_{1}~,
\end{equation}

where we take the mean absolute error between the generator network's predicted value~$\hat{y}_i$ and the ground-truth value~$y_i$, respectively, over every pixel $i$.

\subsection{Adversarial Loss}

Next, we define an adversarial loss to encourage the production of high-resolution images~$\hat{y}$ closer to the manifold of natural-looking digital-elevation-model images.
To do so, we introduce the standard discriminator in the form of $D(y) = \sigma(C(y))$, where~$\sigma$ is the sigmoid activation function and $C(y)$ is the raw, non-transformed output from a discriminator neural network acting on high-resolution image~$y$.
The ESRGAN model \citep{WangESRGANEnhancedSuperResolution2019}, however, employs an improved relativistic-average discriminator \citep{Jolicoeur-Martineaurelativisticdiscriminatorkey2018} denoted by~$D_{\text{Ra}}$.
It is defined as $D_{\text{Ra}}(y,\hat{y}) = \sigma(C(y) - \mathbb{E}_{\hat{y}}[C(\hat{y})])$, where $\mathbb{E}_{\hat{y}}[\cdot]$ is the arithmetic mean operation carried out over every generated image $\hat{y}$ in a mini batch.
We use a binary cross-entropy loss as the discriminator's loss function defined as follows:

\begin{equation}\label{eq:A2}
  L_{\mathrm{D}}^{\text{Ra}} = - \mathbb{E}_y[\ln(D(y,\hat{y}))] - \mathbb{E}_{\hat{y}}[\ln(1 - D(\hat{y},y))].
\end{equation}

The generator network's adversarial loss is in a symmetrical form:

\begin{equation}\label{eq:A3}
  L_{\mathrm{G}}^{\text{Ra}} = - \mathbb{E}_y[\ln(1 - D(y,\hat{y}))] - \mathbb{E}_{\hat{y}}[\ln(D(\hat{y},y))].
\end{equation}

\subsection{Topographic Loss}

We further define a topographic loss so that the elevation values in the super-resolved image make topographic sense with respect to the original low-resolution image.
Specifically, we want the mean value of each 4\,$\times$\,4 grid on the predicted super-resolution (DeepBedMap) image to closely match its spatially corresponding 1\,\unit{pixel}\,$\times$\,1\,\unit{pixel} area on the low-resolution (BEDMAP2) image.

First, we apply a 4\,$\times$\,4 mean pooling operation on the generator network's predicted super-resolution image:

\begin{equation}\label{eq:A4}
  \bar{\hat{y}}_j = \dfrac{1}{n} \sum\limits_{i=1}^n \hat{y}_i~,
\end{equation}

where $\bar{\hat{y}}_j$ is the mean of all predicted values~$\hat{y}_i$ across the 16 super-resolved pixels~$i$ within a 4\,$\times$\,4 grid corresponding to the spatial location of 1 low-resolution pixel at position $j$.
Following this, we can compute the topographic loss as follows:

\begin{equation}\label{eq:A5}
  L_{\mathrm{T}} = \dfrac{1}{m} \sum\limits_{i=1}^m ||\bar{\hat{y}}_j - x_j||_{1}~,
\end{equation}

where we take the mean absolute error between the mean of the 4\,$\times$\,4 super-resolved pixels calculated in Eq.~(\ref{eq:A4}) $\bar{\hat{y}}_j$ and those of the spatially corresponding low-resolution pixel~$x_j$, respectively, over every low-resolution pixel~$j$.

\subsection{Structural Loss}

Lastly, we define a structural loss that takes into account luminance, contrast and structural information between the predicted and ground-truth images.
This is based on the structural similarity index \citep[SSIM;][]{WangImageQualityAssessment2004} and is calculated over a single window patch as

\begin{equation}\label{eq:A6}
  \text{SSIM} (\hat{y}, y) = \dfrac{(2\mu_{\hat{y}}\mu_y + c_1)(2\sigma_{{\hat{y}}y} + c_2)}{(\mu_{\hat{y}}^2 + \mu_y^2 + c_1)(\sigma_{\hat{y}}^2 + \sigma_y^2 + c_2)},
\end{equation}

where $\mu_{\hat{y}}$ and $\mu_y$ are the arithmetic mean of predicted image~${\hat{y}}$ and ground-truth image~$y$, respectively, over a single window that we set to 9~pixels\,$\times$\,9~pixels; $\sigma_{{\hat{y}}y}$ is the covariance of~${\hat{y}}$ and~$y$; $\sigma_{\hat{y}}^2$ and $\sigma_y^2$ are the variances of~${\hat{y}}$ and~$y$, respectively; and~$c_1$ and~$c_2$ are two variables set to $0.01^2$ and $0.03^2$ to stabilize division with a weak denominator.
Thus, we can formulate the structural loss as follows:

\begin{equation}\label{eq:A7}
  L_{\mathrm{S}} = 1 - \dfrac{1}{p} \sum\limits_{i=1}^p \text{SSIM} (\hat{y}, y)_p~,
\end{equation}

where we take $1$ minus the mean of all structural similarity values $\text{SSIM}(\hat{y}, y)$ calculated over every patch~$p$ obtained via a sliding window over the predicted image~${\hat{y}}$ and ground-truth image~$y$.

\subsection{Total Loss Function}

Finally, we compile the loss functions for the discriminator and generator networks as follows:

\begin{align}
  & L_{\mathrm{D}} = L_{\mathrm{D}}^{\text{Ra}}, \label{eq:A8}\\
  & L_{\mathrm{G}} = \eta L_1 + \lambda L_{\mathrm{G}}^{\text{Ra}} + \theta L_{\mathrm{T}} + \zeta L_{\mathrm{S}}~, \label{eq:A9}
\end{align}

where $\eta$, $\lambda$, $\theta$ and $\zeta$ are the scaled weights for the content $L_1$, adversarial $L_{\mathrm{D}}$, topographic $L_{\mathrm{T}}$ and structural losses $L_{\mathrm{S}}$, respectively (see Table~\ref{table:B1} for values used).
The loss functions $L_{\mathrm{D}}$ and $L_{\mathrm{G}}$ are minimized in an alternate $1:1$ manner so as to solve the entire generative adversarial network's objective function defined in Eq.~(\ref{eq:4}).


\section{Neural Network Training Details} \label{appendix:B}

The neural networks were developed using Chainer v7.0.0 \citep{TokuiChainerDeepLearning2019} and trained using full-precision (floating point 32) arithmetic.
Experiments were carried out on four graphical processing units (GPUs), specifically two Tesla P100 GPUs and two Tesla V100 GPUs.
On the Tesla V100 GPU setup, one training run with about 150 epochs takes about 30\,min.
This is using a batch size of 128 on a total of 3826 training image tiles, with 202 tiles reserved for validation, i.e. a $95/5$ training/validation split.
We next describe the method used to evaluate each DeepBedMap candidate model, as well as the high-level way in which we semi-automatically arrived at a good model via semi-automatic hyperparameter tuning.

%TB1
\begin{table}[b!]
  \caption{Optimized hyperparameter settings.}
  \label{table:B1}
  \begin{tabularx}{\textwidth}{lrr}
  \tophline
  Hyperparameter                                       & Setting             & Tuning range                           \\
  \middlehline
  Learning rate (for both generator and discriminator) & $1.7\times 10^{-4}$ & $2\times 10^{-4}$ to $1\times 10^{-4}$ \\
  Number of residual-in-residual blocks                & 12                  & 8 to 14                                \\
  Mini-batch size                                      & 128                 & 64 or 128                              \\
  Number of epochs                                     & 140                 & 90 to 150                              \\
  Residual scaling                                     & 0.2                 & 0.1 to 0.5                             \\
  Content-loss weighting $\eta$                        & $1\times 10^{-2}$   & Fixed                                  \\
  Adversarial-loss weighting $\lambda$                 & $2\times 10^{-2}$   & Fixed                                  \\
  Topographic-loss weighting $\theta$                  & $2\times 10^{-3}$   & Fixed                                  \\
  Structural-loss weighting $\zeta$                    & 5.25                & Fixed                                  \\
  He normal initialization scaling                     & 0.1                 & Fixed                                  \\
  Adam optimizer epsilon                               & 0.1                 & Fixed                                  \\
  Adam optimizer beta1                                 & 0.9                 & Fixed                                  \\
  Adam optimizer beta2                                 & 0.99                & Fixed                                  \\
  \bottomhline
  \end{tabularx}
\end{table}

To check for overfitting, we evaluate the generative adversarial network model using the validation dataset after each epoch using two performance metrics~-- a peak signal-to-noise ratio (PSNR) metric for the generator and an accuracy metric for the discriminator.
Training stops when these validation performance metrics show little improvement, roughly at 140~epochs.

Next, we conduct a full evaluation on an independent test dataset, comparing the model's predicted grid output with actual ground-truth \textit{xyz} points.
Using the ``grdtrack'' function in Generic Mapping Tools v6.0 \citep{WesselGenericMappingTools2019}, we obtain the grid elevation at each ground-truth point and use it to calculate the elevation error on a point-to-point basis.
All of these elevation errors are then used to compute a root mean square error (RMSE) statistic over this independent test site.
This RMSE value is used to judge the model's performance in relation to baseline bicubic interpolation and is also the metric minimized by a hyperparameter optimization algorithm which we will describe next.

Neural networks contain a lot of hyperparameter settings that need to be decided upon, and generative adversarial networks are particularly sensitive to different hyperparameter settings.
To stabilize model training and obtain better performance, we tune the hyperparameters (see Table~\ref{table:B1}) using a Bayesian approach.
Specifically, we employ the Tree-structured Parzen Estimator \citep{BergstraAlgorithmsHyperparameterOptimization2011} from the Optuna v2.0.0 \citep{AkibaOptunaNextgenerationHyperparameter2019} library with default settings as per the Hyperopt library \citep{BergstraHyperoptPythonlibrary2015}.
Given that we have four GPUs, we choose to parallelize the hyperparameter tuning experiments asynchronously between all four devices.
The estimator first conducts 20 random experimental trials to scan the hyperparameter space, gradually narrowing down its range to a few candidate hyperparameters in subsequent experiments.
We set each GPU to run a target of 60 experimental trials (i.e. a total of 240), though unpromising trials that have exploding or vanishing gradients are pruned prematurely using the Hyperband algorithm \citep{LiHyperbandNovelBanditBased2018} to save on time and computational resources.
The top models from these experiments undergo further visual evaluation, and we continue to conduct further experiments until a suitable candidate model is found.

% \noappendix       %% use this to mark the end of the appendix section. Otherwise the figures might be numbered incorrectly (e.g. 10 instead of 1).
\end{appendices}




%% The following commands are for the statements about the availability of data sets and/or software code corresponding to the manuscript.
%% It is strongly recommended to make use of these sections in case data sets and/or software code have been part of your research the article is based on.

\section{Code and data availability}

Python code for data preparation, neural network model training and visualization of model outputs is freely available at \url{https://github.com/weiji14/deepbedmap} (last access: 9~August~2020) and at \url{https://doi.org/10.5281/zenodo.3752613} \citep{LeongDeepBedMap2020}.
Neural network model training experiment runs are also recorded at \url{https://www.comet.ml/weiji14/deepbedmap} (last access: 9~August~2020).

The DeepBedMap\_DEM is available from Zenodo at \url{https://doi.org/10.5281/zenodo.3752613}  \citep{LeongDeepBedMap2020}.
The Pine Island Glacier dataset \citep{BinghamDiverselandscapesPine2017} is available on request from Robert Bingham.
The Carlson Inlet dataset \citep{KingIcestreamnot2011} is available on request from Edward King.
Bed elevation datasets from Wilkes Subglacial Basin \citep{FerraccioliAirborneradarbed2018} and Rutford Ice Stream \citep{KingSubglaciallandformsRutford2016} are available from the British Antarctic Survey's Polar Data Centre (\url{https://ramadda.data.bas.ac.uk}, last access: 14~January~2020).
Other Antarctic bed elevation datasets are available from the Center for Remote Sensing of Ice Sheets (\url{https://data.cresis.ku.edu/data/rds}, last access: 15~August~2019) or from the National Snow and Ice Data Center (\url{https://nsidc.org/data/IRMCR2/versions/1}, last access: 15~August~2019).
BEDMAP2 \citep{FretwellBedmap2improvedice2013} and REMA \citep{HowatReferenceElevationModel2018} are available from the Polar Geospatial Center (\url{http://data.pgc.umn.edu}, last access: 30~August~2019).
MEaSUREs Ice Velocity data \citep{MouginotMEaSUREsPhaseMap2019} are available from NSIDC (\url{https://nsidc.org/data/nsidc-0754/versions/1}, last access: 31~August~2019).
Antarctic snow accumulation data \citep{ArthernAntarcticsnowaccumulation2006} are available from the British Antarctic Survey (\url{https://secure.antarctica.ac.uk/data/bedmap2/resources/Arthern_accumulation}, last access: 17~June~2019).


% \authorcontribution{
%   WJL was responsible for data curation, formal analysis, methodology, software, visualization and writing the original draft.
%   HJH was responsible for funding acquisition and supervision.
%   Both authors conceptualized the work and contributed to the reviewing and editing stages of the writing.
% } %% this section is mandatory

\section{Acknowledgements}

We are grateful to Robert Bingham and Edward King for the Pine Island Glacier and Carlson Inlet data and to all the other researchers in the British Antarctic Survey and Operation IceBridge team for providing free access to the high-resolution bed elevation datasets around Antarctica.
A special thanks to Ruzica Dadic for her help in reviewing draft versions of this paper.
This research was funded by the Royal Society of New Zealand's Rutherford Discovery Fellowship (contract RDF-VUW1602), with additional support from the Erasmus+ programme and International Glaciological Society early-career travel award for presenting earlier versions of this work at the 2019 EGU General Assembly and IGS Symposium on Five Decades of Radioglaciology.

Chapter 2, in full, is a reprint of material as it appears in The Cryosphere:
Leong, W. J., \& Horgan, H. J. (2020). DeepBedMap: A deep neural network for resolving the bed topography of Antarctica. The Cryosphere, 14(11), 3687–3705. \url{https://doi.org/10.5194/tc-14-3687-2020}.
The paper was edited by Olivier~Gagliardini and reviewed by Martin~Siegert and one anonymous referee.


\chapter{The role of subglacial topography on Antarctic ice flow}
\label{ch:3}
\chaptermark{Basal inversion}
% The role of subglacial topography on Antarctic ice flow

\section*{Abstract}

To examine the influence of bed topography on ice motion, we conduct an inversion experiment using two beds with different spatial resolutions - the $\SI{250}{\metre}$ resolution DeepBedMap\_DEM and the $\SI{500}{\metre}$ resolution BedMachine Antarctica.
At the bed of a glacier, form drag and skin drag resists the driving stress of the overriding ice.
We examine how an increase in form drag due to higher resolution topography affects the remaining skin drag component which is controlled by subglacial hydrology and bed material properties.
An iterative least-squares control method basal inversion is performed using the Ice-sheet and Sea-level System Model (\gls{ISSM}) with L1L2 stress balance equations and a Schoof-type Coulomb-limited sliding law.
Comparison of inverted fields between DeepBedMap\_DEM and BedMachine Antarctica do not appear to be significantly different.
Basal drag and slipperiness fields are slightly higher, while effective pressure is generally lower when using the higher-resolution DeepBedMap\_DEM, due to a higher bed elevation model and thinner ice inducing less overburden pressure.
Previous inversion studies using a Weertman-style sliding relation over Pine Island Glacier indicated that high resolution bed topography (i.e. more form drag) reduced basal drag (i.e. less skin drag), a conclusion that is not supported in this work as no noticeable decrease in basal drag (i.e. skin drag) was observed when higher resolution topography (i.e. more form drag) was used.
These findings warrant more investigation into the dependence of ice sheet modelling on higher spatial resolution bed topographies ($<= \SI{100}{\metre}$) which increases form drag, and time-dependent subglacial hydrology which influences skin drag, over more diverse subglacial settings in West Antarctica.


\section{Introduction}

The bed topography of Antarctica is a critical parameter in ice sheet modelling studies as it affects basal slip - the combined motion of ice from sliding and bed deformation \citep[][p.223]{Cuffeyphysicsglaciers2010}.
Understanding the slip relation remains one of the biggest uncertainties in our determination of the Antarctic ice sheet's contribution to sea level rise \citep[e.g.][]{RitzPotentialsealevelrise2015,BulthuisUncertaintyquantificationmulticentennial2019}.
The basal conditions of an ice sheet vary spatially, and are often determined via inverse methods.
Inverse methods allow us to use ice surface observations coupled with ice physics to solve for critical basal parameters needed for ice sheet modelling \citep[e.g.][]{MacAyealbasalstressdistribution1992,JoughinBasalconditionsPine2009,MorlighemInversionbasalfriction2013}, and also to resolve bed topography \citep[e.g.][]{MorlighemDeepglacialtroughs2019,LeongDeepBedMapdeepneural2020}.
These indirect empirical methods are used because sampling subglacial sediments under the Antarctic ice sheet is logistically challenging \citep[e.g.][]{Siegertassessmentdeephotwater2014,TulaczykWISSARDSubglacialLake2014}, though having direct access to sediments would allow us to use physically based methods of quantifying basal slip \citep[e.g.][]{Zoetsliplawglaciers2020}.

This chapter explores the interplay between two components of frictional drag at the ice-bed interface - form drag and skin drag \citep{SchoofBasalperturbationsice2002,BinghamDiverselandscapesPine2017,Kyrke-SmithRelevanceDetailBasal2018,Minchewuniversalglacierslip2020}.
Form drag occurs from long-wavelength bed topography opposing ice motion \citep[][]{WeertmanSlidingGlaciers1957,SchoofBasalperturbationsice2002}, while skin drag arises in relation to lubricating water and bed material properties \citep{IversonExperimentsdynamicssedimentary2015}.
To study this, we look at how a high resolution (\SI{250}{\metre}) rougher bed \citep{LeongDeepBedMap2020} compares with a medium resolution (\SI{500}{\metre}) smoother bed \citep{MorlighemMEaSUREsBedMachineAntarctica2020}.
The spatially varying basal drag of the bed can be worked out by using bed topography to calculate form drag \citep{SchoofBasalperturbationsice2002}, and then taking the remainder balance as skin drag which is determined by bed material properties and subglacial hydrology \citep{Kyrke-SmithRelevanceDetailBasal2018}. % TODO define how form/skin drag is partitioned
% A secondary objective is
In the following, we review two key parameters which ice sheet models are sensitive to - multi-scale bed topography features and sliding law relations.



% \subsection{Sensitivity of ice sheet models to basal parameters}

% The physical basis of basal slip over deformable till is provided by Equation \eqref{eq:3.1},~\eqref{eq:3.2},
% The effect of subglacial water on ice flow is likely geographically dependent, in that different basal conditions apply in different parts of Antarctica.

\subsection{Bed Topography}

\citet{DurandImpactbedrockdescription2011} performed a sensitivity analysis on different spatial resolution beds using a 2-D flowline model and suggested that a spatial resolution of at least \SI{1}{\kilo\metre} was needed to accurately model coastal outlet glaciers.
This was followed by the creation of BEDMAP2 \citep{FretwellBedmap2improvedice2013} - a \SI{1}{\kilo\metre} resolution Antarctic bed elevation model resampled from a \SI{5}{\kilo\metre} grid, which became the de facto standard for ice sheet models \citep[see e.g.][]{SeroussiinitMIPAntarcticaicesheet2019}.
However, there have been issues reported in relation to BEDMAP2's ice thickness at the former grounding line of Pine Island Glacier that did not respect mass conservation \citep{RignotWidespreadrapidgrounding2014}.
High resolution bed elevations are especially needed close to grounding zones \citep[e.g.][]{SchoofIcesheetgrounding2007,CornfordAdaptivemeshrefinement2016}, and there is a need for high resolution ($<= \SI{500}{\metre}$) bed elevation models that preserve topographic details, so that form drag - resistance to ice flow due to bed topography - is correctly accounted for \citep{BinghamDiverselandscapesPine2017,Kyrke-SmithRelevanceDetailBasal2018}.
The limitations of BEDMAP2 motivated the creation of higher resolution ($<=$ \SI{500}{\metre}) bed \gls{DEM}s with two important attributes - large-scale topography that is consistent with mass conservation requirements \citep[e.g.][]{Morlighemmassconservationapproach2011}, and preservation of fine-scale roughness details \citep[e.g.][]{GoffConditionalsimulationThwaites2014,Grahamhighresolutionsyntheticbed2017}, both of which are needed to fully capture ice-sheet dynamics.

\paragraph{Long-wavelength details captured by mass conservation}

Ice sheet models include the physical laws of mass conservation, and this can be used to invert ice surface observations into bed elevation \citep[e.g.][]{Morlighemmassconservationapproach2011}.
Using an ensemble of 2-D shelfy stream approximation \citep[SSA;][]{MacAyealLargescaleiceflow1989} models, \citet{SchlegelExplorationAntarcticIce2018} performed a comparison of bedrock topography over the Antarctic ice sheet, and their informed bound (using regionally based boundary conditions) experiments showed that using an improved BEDMAP2 mass conserving grid \citep[cf.][]{RignotWidespreadrapidgrounding2014} produced $0.3292\pm\SI{0.1019}{\metre}$ \gls{SLE}, while the older BEDMAP \citep{LytheBEDMAPnewice2001} grid produced a $0.1845\pm\SI{0.08834}{\metre}$ \gls{SLE} over a \SI{100}{\year} period.
This represents a mean increase of \SI{0.14}{m} \gls{SLE} from the improved resolution BEDMAP2 mass conserving grid, but also an 18\% increase in the \gls{SLE} uncertainty which suggests a wider range of grounding line retreat scenarios with implications for ice sheet stability.
On a catchment scale, \citet{NiasNewMassConserving2018} produced a mass-conserving bed topography over Pine Island Glacier and found that modelled sea level contributions over a \SI{50}{\year} period was consistently higher using the mass-conserving bed (\SIrange{0.08}{0.10}{\milli\metre\per\year} \gls{SLE}) than with BEDMAP2 (\SIrange{0.04}{0.07}{\milli\metre\per\year} \gls{SLE}) irrespective of the sliding law used.
Since then, this mass conservation technique has been applied to other key outlet glaciers and ice streams in BedMachine Antarctica \citep{MorlighemMEaSUREsBedMachineAntarctica2020}, identifying new areas vulnerable to Marine Ice Sheet Instability such as at Ninnis Glacier in George V Land and Denman Glacier in East Antarctica which both rest on retrograde slopes \citep{MorlighemDeepglacialtroughs2019}.
This medium resolution (\SI{500}{\metre}) mass conserving BedMachine Antarctica grid \citep[version 2,][]{MorlighemMEaSUREsBedMachineAntarctica2020} will be used in our basal inversion experiment.

\paragraph{Short-wavelength details from spatial statistics}

% the super-resolution DeepBedMap\_DEM \citep{LeongDeepBedMapdeepneural2020}
Using an ensemble of L1L2 \citep{Hindmarshnumericalcomparisonapproximations2004,SchoofThinFilmFlowsWall2010} ice sheet models, \citet{SunDynamicresponseAntarctic2014} showed that the phase of low-frequency noise ($>\SI{10}{\kilo\metre}$) was more important than the phase of high-frequency noise ($<\SI{1}{\kilo\metre}$) of the same amplitude in their model projection runs.
Still, the ice sheet model was sensitive to the amplitude of the short-wavelength/high-frequency noise even if the phase of that noise was not so important.
Spatial statistical methods have since been used to retrieve this topographical noise along radio-echo sounding (\gls{RES}) transects and apply it spatially to a larger regional extent.
Over Thwaites Glacier, \citet{GoffConditionalsimulationThwaites2014} captured the different roughness statistics between highland and lowland parts, and applied these to the whole catchment using a conditional simulation method which allowed fine-scale roughness and channelized morphologies to be resolved better than ordinary kriging.
\citet{Grahamhighresolutionsyntheticbed2017} generated a synthetic high-resolution ($\SI{100}{\metre}$) grid using a two-step approach, with the low-frequency component derived from BEDMAP2, and a high-frequency component derived from radar point data from the Investigating the Cryospheric Evolution of the Central Antarctic Plate (ICECAP) and BEDMAP compilation.
\citet{LeongDeepBedMapsuperresolutiondeep2019} used 2-D gridded \gls{RES} ground-truth data with ice surface observations of elevation, velocity and snow accumulation to train a super-resolution deep neural network model, using it to produce a \SI{250}{\metre} resolution grid of the Antarctic ice sheet.
This super-resolution (\SI{250}{\metre}) DeepBedMap\_DEM grid \citep[version 1.1,][]{LeongDeepBedMap2020} will be used in our basal inversion experiment.



\subsection{Basal drag}

% \subsection{Form drag vs Skin drag}

Form drag arises due to large-scale bed topography which opposes the motion of ice as it flows and deforms around obstacles on a rigid bed \citep[][]{WeertmanSlidingGlaciers1957,SchoofBasalperturbationsice2002}.
Skin drag occurs in relation to bed material properties and in the presence of water that acts to lubricate the ice-bed interface and allow subglacial till to deform more easily \citep{IversonExperimentsdynamicssedimentary2015}.
The physical mechanisms of these two processes may differ, but the mathematical formulation and parameterizations for both are similar \citep[Fig.~\ref{fig:1.3},][]{Minchewuniversalglacierslip2020}.

\subsubsection{Physical basis of basal slip}

A basal slip law that relates basal drag or basal shear stress $\tau_b$ with effective normal stress (pressure) $N$ and basal slip velocity $u_b$ is given by \citet[][eq.1]{Zoetsliplawglaciers2020}:

\begin{equation}
  \tau_b = \min [ N \tan(\phi), (C u_b)^{1/m} ] \label{eq:3.1}
\end{equation}

where $\phi$ is the friction angle of the basal till, such that $\tan(\phi)$ is the friction coefficient, $m$ is a slip law exponent, and $C$ is a constant dependent on a transitional velocity $u_t$ and bed roughness, formulated as $C = \frac{(N \tan(\phi))^{m}}{u_t}$.
The approximate continuous form of Equation \eqref{eq:3.1} \citep[][eq.3]{Zoetsliplawglaciers2020} is as follows:

\begin{equation}
  \tau_b = N \tan(\phi) \left( \frac{u_b}{u_b + u_t} \right)^{1/p} \label{eq:3.2}
\end{equation}

where $p$ is a slip exponent, experimentally determined to be $\sim5$ \citep{Zoetsliplawglaciers2020}.

At low basal slip speeds ($u_b \to 0$) typically found over rigid, dry beds, viscous Weertman-style behaviour \citep[$(C u_b)^{1/m}$ in Eq.~\ref{eq:3.1},][]{WeertmanSlidingGlaciers1957} dominates total drag $\tau_b$.
At high basal slip speeds ($u_b \to \infty$) found over water-saturated deforming beds, Coulomb-plastic behaviour \citep[$N \tan(\phi)$ in Eq.~\ref{eq:3.1}, e.g.][]{Schoofeffectcavitationglacier2005,JoughinRegularizedCoulombFriction2019} is more influential.
The upper bound of $\tau_b$ is based on the shear strength of the till, as determined from ring-shear laboratory experiments \citep{Zoetsliplawglaciers2020}.
Past the transition speed $u_t$, basal drag $\tau_b$ becomes independent of basal slip velocity $u_b$ \citep[c.f.][]{StearnsFrictionbeddoes2018}.
The physical formulation of Equation \eqref{eq:3.1} is similar to that of other Coulomb-based parametrizations \citep[e.g.][]{TsaiMarineicesheetprofiles2015,JoughinRegularizedCoulombFriction2019}, but while previous implementations fixed $C$ to estimate $u_t$ via inversion (\textit{a posteriori}), Equation \eqref{eq:3.1} is formulated on a physical basis with $u_t$ derived independently \citep[using clast size and till placement, see eq.2 in][]{Zoetsliplawglaciers2020} and then used to directly determine $C$ (\textit{a priori}).


\subsubsection{Sliding laws - linear viscous to Coulomb-limited}

% TODO throw in Coulomb sliding law review from \citet{NobleSensitivityAntarcticIce2020}

It is not feasible to access subglacial till across a glaciated area to physically determine $u_t$ using Equation \eqref{eq:3.1}, so parameterized basal slip relations are used in practice.
Sensitivity analyses conducted in multiple ice sheet model comparison studies \citep[e.g.][]{SeroussiinitMIPAntarcticaicesheet2019,SunAntarcticicesheet2020,ZhangcomparisontwoStokes2017} have explored the use of different basal slip relations or sliding laws (see Fig.~\ref{fig:sliding_laws}), and the different tunable parameters that come with them.
\citet{Gillet-ChauletAssimilationsurfacevelocities2016} used surface velocity observations from 1996 to 2010 over Pine Island Glacier and suggested that a non-linear Weertman stress exponent of $m >= 5$ best matched observed flow accelerations, with the assumption that the basal slipperiness coefficient $C$ remained constant over the period.
In the Antarctic BUttressing Model Intercomparison Project \citep[ABUMIP;][]{SunAntarcticicesheet2020}, the ice-shelf removal or `float-kill' experiment (ABUK) noted how ice sheet models implementing linear Weertman/Coulomb friction laws tend to show lower ice loss (\SI{3.07}{\metre} \gls{SLE}) than models which implemented pseudo-plastic (\SI{4.41}{\metre} \gls{SLE}) or plastic (\SI{10.20}{\metre} \gls{SLE}) sliding laws.
Ice sheet models implementing such non-linear/plastic sliding laws generally lead to faster grounding line retreat than those using viscous linear sliding laws over the Amundsen Sea Embayment area \citep[e.g.][]{JoughinBasalconditionsPine2009,RitzPotentialsealevelrise2015,BrondexSensitivitygroundingline2017,BulthuisUncertaintyquantificationmulticentennial2019}.

\begin{figure}[htbp]
  \includegraphics[width=1.0\textwidth]{chp3_fig_sliding_law_comparisons.jpg}
  \caption[Comparison of Weertman, Budd, Schoof and Tsai sliding laws]{
    Comparison of linear (Weertman, Budd) and non-linear, Coulomb limited (Schoof and Tsai) sliding laws.
    Basal velocity $u_b$ is plotted on the vertical y-axis and effective pressure $N$ is plotted on the horizontal x-axis.
    Coloured iso-lines represent basal drag ($\tau_b$) values ranging from \SI{0.04}{\mega\pascal} to \SI{0.2}{\mega\pascal}.
    Figure is from \citet{BrondexSensitivitygroundingline2017}.
  }
  \label{fig:sliding_laws}
\end{figure}

Other intercomparisons have judged the ability of different sliding laws to properly model grounding line dynamics, with overarching support for effective pressure ($N$) dependent sliding laws such as Budd-type \citep{BuddEmpiricalStudiesIce1979}, Tsai-type \citep{TsaiMarineicesheetprofiles2015} or Schoof-type \citep{Schoofeffectcavitationglacier2005} formulations instead of a classical Weertman-type \citep{WeertmanSlidingGlaciers1957} relation \citep[see Fig.~\ref{fig:sliding_laws},][]{BrondexSensitivitygroundingline2017,BrondexSensitivitycentennialmass2019}.
In particular, \citet{JoughinRegularizedCoulombFriction2019} noted how the effects of cavitation on sliding \citep[see][]{Schoofeffectcavitationglacier2005} appear important at Pine Island Glacier, recommending a regularized Coulomb sliding law with $m = 3$ and velocity threshold $u_o = 300\,\si{\metre\per\year}$ to reliably reproduce glacier behaviour in a manner applicable to both weak till and hard bedrock areas.
Coulomb-limited sliding relations are generally favoured, such as a Schoof-type law for its ability to transition continuously between Weertman and Coulomb sliding regimes in different settings \citep{BrondexSensitivitygroundingline2017,BrondexSensitivitycentennialmass2019,CornfordResultsthirdMarine2020,Zoetsliplawglaciers2020,Minchewuniversalglacierslip2020}.
Thus, we choose to use the Schoof sliding law \citep[adapted from][eq. 6.2]{Schoofeffectcavitationglacier2005}, presented below in a form similar to that of Equation \eqref{eq:3.2}:

\begin{equation}
  \tau_b = N C \left( \frac{\Lambda}{\Lambda + \Lambda_0} \right)^{1/m}, \Lambda = \frac{u_b}{N^m} \label{eq:3.3}
\end{equation}

where $\tau_b$ is basal drag, $N$ is effective pressure, $C$ is the Schoof friction coefficient less than the maximum bed slope (c.f. $\tan(\phi)$ in Eq.~\eqref{eq:3.2}), $u_b$ is basal velocity, $\Lambda_0$ is a maximum threshold that satisfies Iken's bound (c.f. $u_t$ in Eq.~\eqref{eq:3.2}), and $m$ is a power law exponent.
Equation \eqref{eq:3.3} will be reformulated as Equation \eqref{eq:schoof} in the Methods section to follow the \gls{ISSM} \citep{LarourContinentalscalehigh2012} Schoof sliding law implementation.


% TODO!! Zoet & Iverson 2020 and Michew & Joughin 2020
% Use https://science.sciencemag.org/content/sci/368/6486/29/F1.large.jpg in Thesis Introduction!!


\subsection{Previous work} % Relevance of basal roughness

\citet{Kyrke-SmithRelevanceDetailBasal2018} ran a Full Stokes ice sheet model inversion experiment with a Weertman-style sliding law on two different beds - a high-resolution (\SI{50}{\metre}) bed topography from DELORES data \citep{BinghamDiverselandscapesPine2017} and on a low-resolution (\SI{1000}{\metre}) BEDMAP2 \citep{FretwellBedmap2improvedice2013} - to compare the balance of slipperiness ($C$) and basal drag ($\tau_b$) estimates.
They reported that spatial patterns of both slipperiness $C$ and basal drag $\tau_b$ were similar across both high- and low-resolution experiments.
However, the magnitude of the spatially averaged basal drag $\tau_b$ (i.e. skin drag) was consistently lower when using the high-resolution DELORES \gls{DEM} than with the low-resolution BEDMAP2 owing to more form drag being accounted for.
Over the iSTARt1 area for example, the high-resolution DELORES \gls{DEM} had $\bar{\tau_b} = \SI{6.8}{\kilo\pascal}$, low-resolution BEDMAP2 had $\bar{\tau_b} = \SI{9.6}{\kilo\pascal}$, while the flat bed control experiment had $\bar{\tau_b} = \SI{11.9}{\kilo\pascal}$.
The lower mean basal drag $\bar{\tau_b}$ was attributed to an increase in form drag introduced by the high-resolution DELORES DEM that induces more resistive stresses on ice.
The skin drag component must therefore decrease, corresponding to a lower slipperiness $C$, in order for the modelled velocity to match the observed velocity in the inversion process.
The implications of this are that both basal drag $\tau_b$ and slipperiness $C$ may be overestimated for ice sheet models using parameterizations based on low-resolution bed topography models like BEDMAP2 \citep{Kyrke-SmithRelevanceDetailBasal2018}.
Here, we expand on the work of \citet{Kyrke-SmithRelevanceDetailBasal2018} by running an L1L2 ice sheet model \citep[\gls{ISSM};][]{LarourContinentalscalehigh2012} using a Coulomb-limited Schoof-type sliding law \citep{JoughinRegularizedCoulombFriction2019,Schoofeffectcavitationglacier2005} over the larger spatial area of the main trunk of Pine Island Glacier (see Fig.~\ref{fig:topo_and_speed}).

% TODO paraphrase
% It therefore makes sense that the derived skin drag is consistently lower when more detailed bed variations are included in the domain; the form drag transmitted by the bed variations balances more of the driving stress.

\clearpage
\section{Methods} \label{sec:methods}

\subsection{Ice Sheet model set-up}

% Refer to https://tc.copernicus.org/articles/12/3861/2018/ setup
% Also https://doi.org/10.1017/jog.2020.8, ISSM model

\subsubsection{Stress balance model}

We use an L1L2 model \citep{Hindmarshnumericalcomparisonapproximations2004,SchoofThinFilmFlowsWall2010} implemented in the Ice-sheet and Sea-level System Model \citep[\gls{ISSM}, v4.18;][]{LarourContinentalscalehigh2012}.
\gls{ISSM} is a thermomechanical finite-element ice flow model that follows physical laws for the conservation of mass, momentum and energy, coupled with constitutive material laws and boundary conditions \citep{LarourContinentalscalehigh2012}.
Ice is treated as an viscous incompressible material \citep{Cuffeyphysicsglaciers2010} with a non-linear ice effective viscosity $\mu$ following a Norton-Hoff rheology law \citep[Glen's flow law,][]{Glencreeppolycrystallineice1955}:

\begin{equation}
  \mu = \frac{B}{2 \dot{\epsilon}_e^{(n-1)/n}}
\end{equation}

where $B$ is the temperature dependent ice hardness (rigidity), $\dot{\epsilon}_e$ is the effective strain rate tensor, $n$ is Glen's flow law exponent.

\subsubsection{Sliding law}

Within \gls{ISSM}, we applied a Schoof-type basal sliding law \citep{Schoofeffectcavitationglacier2005}:

% from https://icepack.github.io/notebooks/tutorials/04-synthetic-ice-stream/
% \begin{equation}
%   \tau_b = -\frac{\tau_0|u_b|^{\frac{1}{m} - 1}u_b}{\left(u_0^{\frac{1}{m} + 1} + |u_b|^{\frac{1}{m} + 1}\right)^{\frac{1}{m + 1}}}
% \end{equation}

\begin{equation}
  \tau_b = -\frac{C |u_b|^{m-1} u_b}{\left(1 + \left(\frac{C}{C_{\text{max}}N}\right)^{1/m} + |u_b|\right)^m} \label{eq:schoof}
\end{equation}

where $\tau_b$ is basal drag, $C$ is the Schoof friction coefficient (slipperiness), $C_{\text{max}}$ is a Coloumb friction upper bound parameter, $u_b$ is basal ice velocity, $N$ is effective pressure, and $m$ is a positive power law exponent (note that $m$ here is the reciprocal of $m$ in other studies \citep[e.g.][]{TsaiMarineicesheetprofiles2015}, i.e. $m = 1/m_{\text{Tsai}} = 1/3$).

The friction factor $C_{\text{max}}$ imposes an upper limit on $\tau_b$ which must satisfy Iken's bound \citep{IkenEffectSubglacialWater1981,GagliardiniFiniteelementmodelingsubglacial2007}:

\begin{equation}
  \tau_b \leq N\tan\beta
\end{equation}

where $\beta$ is the maximum local up-slope angle between the bed and mean flow direction.
Based on lab experiments \citep{IversonRingshearstudiestill1998}, $C_{\text{max}}$ is set to between $\tan \SI{10}{\degree} = 0.17$ and $\tan \SI{40}{\degree} = 0.84$, with lower values corresponding to clay-rich tills \citep[, pp.266-267]{Cuffeyphysicsglaciers2010}.
The spatial distribution of $C_{max}$ varies across a glacier but no method yet exists to constrain it, so a uniform value of 0.4 is used instead \citep[0.4 and 0.6 was used in][]{BrondexSensitivitycentennialmass2019}.
% The magnitude of $\tau_b$ has an asymptotic behaviour, tending towards $|u|^{\frac{1}{m}}$ at low ice velocities when $|u| << u_o$, but to $\tau_0$ at high ice velocities when $|u| >> u_o$.
The magnitude of $\tau_b$ has an asymptotic behaviour, such that $\tau_b \sim C u_b^m$ as $N$ trends towards infinity corresponding to a Weertman-type sliding regime, and $\tau_b \sim C_{\text{max}} N$ as $N$ trends towards zero corresponding to a Coloumb-type sliding regime \citep{BrondexSensitivitygroundingline2017}.

The effective pressure ($N$) can be approximated using ice overburden pressure minus water pressure:

\begin{equation}
  N = \rho_i g H - P_w
\end{equation}

where $\rho_i$ is ice density, $g$ is gravitational acceleration, $H$ is ice thickness and $P_w$ is water pressure.

% A Budd-type sliding law \citep{BuddEmpiricalStudiesIce1979} is applied with parameters $p=1$, $q=1$, % TODO check parameters

% \begin{equation}
%   \tau_b = -C N^r |v_b|^{s-1} v_b
% \end{equation}

% where $N$ is the effective pressure at the ice base, $C$ is slipperiness (friction coefficient), $v_b$ is basal velocity.
% r and s are defined as

% \begin{equation}
%   r = \frac{q}{p}, s=\frac{1}{p}
% \end{equation}

% where $p$ and $q$ are sliding law exponents.
% The parameter $s$ is similar to $m$ in a Weertman sliding law, and we set $p=0.2$ and $q=0.2$ to obtain an $m=5$ non-linear sliding law.

%TODO use m=8 (not 1) and see what happens

\subsubsection{Boundary conditions} \label{sec:boundary_conditions}

\begin{figure}[htbp]
  \centering
  \includegraphics[width=0.99\textwidth]{figures/chp3_fig_topography_and_speed}
  \caption[Topographies and Speed over Pine Island Glacier]{
    Topographies and speed over Pine Island Glacier study area.
    a) Ice surface elevation from \citet{HowatReferenceElevationModel2019}.
    b) Ice surface speed (contours in white) from \citet{MouginotMEaSUREsPhaseMap2019}.
    c) DeepBedMap\_DEM bed elevation.
    d) BedMachine bed elevation.
    Plotted on an Antarctic Stereographic Projection with a standard latitude of 71°S (EPSG:3031).
  }
  \label{fig:topo_and_speed}
\end{figure}

% Get from the PIG_par.py file
The boundary conditions are kept equal across all experiments, aside from modifications to the bed topography (see Fig.~\ref{fig:topo_and_speed}). % TODO and sliding law parameters?
The model's spatial domain is a triangulated static anisotropic adaptive mesh grid with 10 vertical layers, and a spatial resolution of \SIrange{250}{20000}{\metre} (higher resolution over areas of high ice velocity), located over the main trunk of Pine Island Glacier.
The same mesh is used for both experiments.
A minimum ice thickness of \SI{1}{\metre} is set to avoid numerical instabilities.
A stress-free surface is applied at the ice-atmosphere interface $\boldsymbol{\sigma} \cdot \boldsymbol{n} = 0$ where $\boldsymbol{\sigma}$ is the Cauchy stress tensor and $\boldsymbol{n}$ is the unit outward-pointing normal vector \citep[][eq. 18]{LarourContinentalscalehigh2012}.
Surface mass balance, ice temperature and geothermal heat flux are from ALBMAPv1 \citep{LeBrocqimprovedAntarcticdataset2010}.

\subsection{Inverting for basal drag $\tau_b$}

To determine the spatial distribution of basal drag ($\tau_b$), we run an inversion using remotely sensed values of surface velocity $u_s$ (Fig.~\ref{fig:topo_and_speed}b).
This is accomplished using a least-squares approach involving a control method \citep[][eq. 6]{MacAyealbasalstressdistribution1992} which avoids the errors in $\tau_b$ caused by using a direct algebraic inversion approach \citep[e.g.][eq. 4, 5]{MacAyealbasalstressdistribution1992} which is more sensitive to un-physical interpolations of $u_s$.
Following \citet[][eq. 9]{MorlighemInversionbasalfriction2013}, the objective function or cost function $J$ used to minimize the difference in modelled and observed surface velocities is formulated as follows:

\begin{equation}
  J(\boldsymbol{v}, \alpha) = J_{mis} + \gamma_1 J_{1} + \gamma_t J_{2}
\end{equation}

where $\boldsymbol{v}$ is ice surface velocity, $\alpha$ is the parameter being inverted (i.e. the friction coefficient $\ln(C)$), and $\gamma_1, \gamma_t$ are non-dimensional weighting constants.
The initial friction coefficient was set as $\alpha_0 = \SI{10}{(\pascal\,\year\per\metre)}^{1/2}$ for the model domain.
We set $\gamma_1 = 100$ and $\gamma_t = 10^{-7}$ following \citet{MorlighemInversionbasalfriction2013}.

The first term $J_{mis}$ is an $L^2$ misfit calculated as:

\begin{equation}
  J_{mis} = \frac{1}{2} \int_{\Gamma_s} (v_x - v_x^{obs})^2 + (v_y - v_y^{obs})^2 \,d\Gamma_s
\end{equation}

where $v_x$ and $v_y$ are modelled ice surface velocities (in the $x$ and $y$ directions), $v_x^{obs}$ and $v_y^{obs}$ are observed ice surface velocities from \citet{MouginotContinentWideInterferometric2019}, and $\Gamma_s$ is the ice surface domain.

The second term $J_1$ takes the squared natural logarithmic difference between $u$ and $u_{obs}$:

\begin{equation}
  J_1 = \frac{1}{2} \int_{\Gamma_s} \ln \left( \frac{\sqrt{u^2 + v^2} + \epsilon}{\sqrt{u_{obs}^2 + v_{obs}^2} + \epsilon}  \right)^2 \,d\Gamma_s
\end{equation}

where $\epsilon$ is a minimal velocity used to avoid division by zero. This $J_1$ term accounts for the order of magnitude difference in ice flow speed that ranges from $>\SI{1000}{\metre\per\year}$ at ice streams to $<\SI{1}{\metre\per\year}$ in the slow moving interior \citep{MouginotContinentWideInterferometric2019}.

The third term $J_2$ is a Tikhonov regularization term which applies a degree of smoothness to the inversion:

\begin{equation}
  J_2 = \frac{1}{2} \int_{\Gamma_b} \nabla\alpha \cdot \nabla\alpha \,d\Gamma_b
\end{equation}

where $\Gamma_b$ is the ice bed domain, and recall that $\alpha$ is the parameter being inverted.

The inversion is computed by iteratively minimizing the total cost function $J$ until a convergence stopping criterion is reached.
The stopping criterion is set empirically as when the cost function output value reduces by $<0.1$ between successive iterations.
Further details of the inversion process can be found in \citet{MorlighemInversionbasalfriction2013}.

\clearpage
\section{Results} \label{sec:results}

We now present our results over Pine Island Glacier, comparing the medium spatial resolution (\SI{500}{\metre}) smooth bed \citep[BedMachine v2;][]{MorlighemMEaSUREsBedMachineAntarctica2020} with a high spatial resolution (\SI{250}{\metre}) rough bed \citep[DeepBedMap\_DEM v1.1;][]{LeongDeepBedMap2020}.
In the following, we present the inverted spatial distributions of each bed's modelled velocity ($u_b$, Fig.~\ref{fig:velocity}) along with effective pressure ($N$, Fig.~\ref{fig:effective_pressure}), slipperiness ($C$, Fig.~\ref{fig:slipperiness}), and basal drag ($\tau_b$, Fig.~\ref{fig:basal_drag}) fields.
For all the fields above, we also present a transect plot taken along the main trunk of Pine Island Glacier (Fig.~\ref{fig:transect_AB}) and across flow (Fig.~\ref{fig:transect_CD}).

\paragraph{Basal velocity ($u_b$)}

In Fig.~\ref{fig:velocity}, velocity $u_b$ is shown to increase from the upstream part of the glacier (A) to the downstream part of the glacier (B) near the grounding zone for both DeepBedMap and BedMachine.
The mean velocity ($\bar{u_b}$) for DeepBedMap\_DEM (\SI{614}{\metre\per\year}) is slightly higher than that of BedMachine (\SI{552}{\metre\per\year}), a difference of \SI{62}{\metre\per\year}.
Spatially, the difference in velocity between the two beds appears to be greatest near the grounding zone around point B (see also Fig.~\ref{fig:transect_CD}b).

\paragraph{Effective Pressure ($N$)}

In Fig.~\ref{fig:effective_pressure}, effective pressure $N$ is shown to decrease from the upstream part of the glacier (A) to the downstream part of the glacier (B) near the grounding zone for both DeepBedMap and BedMachine.
The mean effective pressure ($\bar{N}$) for DeepBedMap\_DEM ($\SI{1.3389e7}{\pascal}$) is lower than that of BedMachine ($\SI{1.3807e7}{\pascal}$), a difference of $\SI{4.1762e5}{\pascal}$.
The along-flow transect plot (see Fig.~\ref{fig:transect_AB}c) shows that DeepBedMap\_DEM's effective pressure field over the grounding zone (near point B) is noisier than that of BedMachine.

\begin{landscape}
\begin{figure}[htbp]
  \includegraphics[scale=0.65]{figures/chp3_fig4_inverted_bed_velocity}
  \caption[Comparison of basal velocity over Pine Island Glacier for DeepBedMap and BedMachine]{
    Comparison of basal velocity $u_b$ over Pine Island Glacier for DeepBedMap and BedMachine.
    Ice flow direction is from top (point A) to bottom (point B).
    Transect lines cut along flow from upstream point A to downstream point B, and transversely from high elevation point C to low elevation point D.
    Plotted on an Antarctic Stereographic Projection with a standard latitude of 71°S (EPSG:3031).
  }
  \label{fig:velocity}
\end{figure}
\end{landscape}

\begin{landscape}
\begin{figure}[htbp]
  \includegraphics[scale=0.65]{figures/chp3_fig3_inverted_bed_pressure}
  \caption[Comparison of effective pressure over Pine Island Glacier for DeepBedMap and BedMachine]{
    Comparison of effective pressure $N$ over Pine Island Glacier for DeepBedMap and BedMachine.
    Ice flow direction is from top (point A) to bottom (point B).
    Transect lines cut along flow from upstream point A to downstream point B, and transversely from high elevation point C to low elevation point D.
    Plotted on an Antarctic Stereographic Projection with a standard latitude of 71°S (EPSG:3031).
  }
  \label{fig:effective_pressure}
\end{figure}
\end{landscape}

\paragraph{Slipperiness ($C$)}

In Fig.~\ref{fig:slipperiness}, slipperiness $C$ (friction coefficient) is shown to increase from the upstream part of the glacier (A) to the downstream part of the glacier (B) near the grounding zone for both DeepBedMap and BedMachine.
The mean slipperiness ($\bar{C}$) for DeepBedMap\_DEM ($\SI{31.39}{(\pascal\,\year\per\metre)}^{1/2}$) is slightly higher than that of BedMachine ($\SI{31.33}{(\pascal\,\year\per\metre)}^{1/2}$), a difference of $\SI{0.06}{(\pascal\,\year\per\metre)}^{1/2}$.
The along-flow transect plot (see Fig.~\ref{fig:transect_AB}d) shows that DeepBedMap\_DEM is missing one slipperiness peak at Y:-130000 compared to BedMachine Antarctica, owing to an offset of one the ridges.
% Spatially, the difference in slipperiness between the two beds is more pronounced near the shear margins of the ice stream (see across-track transect at Fig.~\ref{fig:transect_CD}d).

\paragraph{Basal drag ($\tau_b$)}

In Fig.~\ref{fig:basal_drag}, basal drag $\tau_b$ is shown to increase from the upstream part of the glacier (A) to the downstream part of the glacier (B) near the grounding zone for both DeepBedMap and BedMachine.
The mean basal drag ($\bar{\tau_b}$) for DeepBedMap\_DEM (\SI{203}{\pascal}) is slightly higher than that of BedMachine (\SI{192}{\pascal}), a difference of \SI{11}{\pascal}.
Spatially, the difference in basal drag between the two beds appears more pronounced near the shear margins along the main trunk of the glacier (see across-track transect at Fig.~\ref{fig:transect_CD}e).

% are more pronounced at the grounding line near point B, with BedMachine showing a peak value of \SI{1.2e14}{\pascal} and DeepBedMap\_DEM showing \SI{1.6e14}{\pascal}, or a difference of \SI{4e13}{\pascal}.

\paragraph{Transect along Pine Island Glacier trunk}

Transect plots are shown along the main trunk of Pine Island Glacier from upstream point A to downstream point B (Fig.~\ref{fig:transect_AB}), and across flow from high elevation point C to low elevation point D (Fig.~\ref{fig:transect_CD}), with DeepBedMap\_DEM \citep{LeongDeepBedMap2020} and BedMachine \citep{MorlighemMEaSUREsBedMachineAntarctica2020} both showing the same broad trends.
Modelled basal velocity ($u_b$) for BedMachine appears to fit closer to the ground-truth velocities observed by remote sensing methods \citep{MouginotMEaSUREsPhaseMap2019} than for DeepBedMap\_DEM (Fig.~\ref{fig:transect_AB}b, Fig.~\ref{fig:transect_CD}b).
Effective pressure ($N$) tends to be lower for DeepBedMap\_DEM than that of BedMachine across both transect lines (Fig.~\ref{fig:transect_AB}c, Fig.~\ref{fig:transect_CD}c), except for near the grounding line close to point B (Fig.~\ref{fig:transect_AB}c).
Slipperiness ($C$, Fig.~\ref{fig:transect_AB}d, Fig.~\ref{fig:transect_CD}d) and basal drag ($\tau_b$, Fig.~\ref{fig:transect_AB}e, Fig.~\ref{fig:transect_CD}e) fields are mostly comparable for both bed topographies.

\begin{landscape}
\begin{figure}[htbp]
  \centering
  \includegraphics[scale=0.65]{figures/chp3_fig2_inverted_bed_slipperiness}
  \caption[Comparison of slipperiness over Pine Island Glacier for DeepBedMap and BedMachine]{
    Comparison of slipperiness $C$ (friction coefficient) over Pine Island Glacier for DeepBedMap and BedMachine.
    Ice flow direction is from top (point A) to bottom (point B).
    Transect lines cut along flow from upstream point A to downstream point B, and transversely from high elevation point C to low elevation point D.
    Plotted on an Antarctic Stereographic Projection with a standard latitude of 71°S (EPSG:3031).
  }
  \label{fig:slipperiness}
\end{figure}
\end{landscape}

\begin{landscape}
\begin{figure}[htbp]
  \includegraphics[scale=0.65]{figures/chp3_fig1_inverted_bed_basal_drag}
  \caption[Comparison of basal drag over Pine Island Glacier for DeepBedMap and BedMachine]{
    Comparison of basal drag $\tau_b$ over Pine Island Glacier for DeepBedMap and BedMachine.
    Ice flow direction is from top (point A) to bottom (point B).
    Transect lines cut along flow from upstream point A to downstream point B, and transversely from high elevation point C to low elevation point D.
    Plotted on an Antarctic Stereographic Projection with a standard latitude of 71°S (EPSG:3031).
  }
  \label{fig:basal_drag}
\end{figure}
\end{landscape}

\begin{figure}[htbp]
  \includegraphics[width=1.0\textwidth]{figures/chp3_fig_transect_plot_AB}
  \caption[Transect plot along Pine Island Glacier trunk]{
    Transect plot along Pine Island Glacier trunk from upstream point A to downstream point B.
    Panels from top to bottom shows comparison of DeepBedMap\_DEM and BedMachine in terms of:
    a) Bed elevation ($z_b$),
    b) Ice surface velocity ($u_s$),
    c) Effective pressure ($N$),
    d) Slipperiness ($C$),
    e) Basal drag ($\tau_b$).
  }
  \label{fig:transect_AB}
\end{figure}

\begin{figure}[htbp]
  \includegraphics[width=1.0\textwidth]{figures/chp3_fig_transect_plot_CD}
  \caption[Transect plot cutting across Pine Island Glacier trunk]{
    Transect plot cutting across Pine Island Glacier trunk from high elevation point C to low elevation point D.
    Panels from top to bottom shows comparison of DeepBedMap\_DEM and BedMachine in terms of:
    a) Bed elevation ($z_b$),
    b) Ice surface velocity ($u_s$),
    c) Effective pressure ($N$),
    d) Slipperiness ($C$),
    e) Basal drag ($\tau_b$).
  }
  \label{fig:transect_CD}
\end{figure}


\clearpage
\section{Discussion and Conclusion}

% \subsection{Basal friction}
The inverted basal drag ($\bar{\tau_b}$, Fig.~\ref{fig:basal_drag}) and slipperiness ($\bar{C}$, Fig.~\ref{fig:slipperiness}) over Pine Island Glacier is slightly higher for DeepBedMap\_DEM than BedMachine Antarctica, while effective pressure ($\bar{N}$, Fig.~\ref{fig:effective_pressure}) is generally lower (see Sec.~\ref{sec:results}).
This mean difference in basal drag from using two different resolution beds is only \SI{11}{\pascal} and not very significant.
In particular, this result is different to the findings of \citet{Kyrke-SmithRelevanceDetailBasal2018} who showed a reduction in basal drag (less skin drag) and slipperiness $C$ when using a higher resolution basal topography (more form drag).
Several factors could account for this discrepancy.

Firstly, the \SI{250}{\metre} difference in spatial resolution between DeepBedMap\_DEM \citep{LeongDeepBedMap2020} and BedMachine Antarctica \citep{MorlighemMEaSUREsBedMachineAntarctica2020} is much smaller than the \SI{950}{\metre} difference between DELORES DEM \citep{BinghamDiverselandscapesPine2017} and BEDMAP2 \citep{FretwellBedmap2improvedice2013}.
This smaller resolution difference may mean that the differences in basal drag (Fig.~\ref{fig:basal_drag}) and slipperiness (Fig.~\ref{fig:slipperiness}) are less apparent.
It could also be a higher resolution ($<= \SI{100}{\metre}$) bed is needed \citep[e.g.][]{Grahamhighresolutionsyntheticbed2017} before the effects of form drag are significant enough to affect slipperiness $C$ and hence basal drag $\tau_b$.
Alternatively, the effects may have been masked when bed elevation is interpolated to the same \SIrange{250}{20000}{\metre} resolution mesh for both experiments (Sec.~\ref{sec:boundary_conditions}), and future studies should investigate such resolution dependence effects.

Secondly, the use of a Schoof-type friction law in our inversion (see \ref{sec:methods}) instead of a Weertman-type sliding law as in \citet{Kyrke-SmithRelevanceDetailBasal2018} may mean that form drag from high resolution topography has less of an effect on basal drag $\tau_b$ (i.e. skin drag).
Following the logic of \citet{Kyrke-SmithRelevanceDetailBasal2018}, a more accurate, high-resolution bed topography model (i.e. more form drag) should lower basal drag ($\tau_b$, i.e. skin drag).
Theoretically, this would shift parts of glaciers flowing at a rate close to the transitional velocity ($u_t$) from a skin drag regime to a transitional or form drag regime \citep[see Fig.~\ref{fig:1.3},][]{Minchewuniversalglacierslip2020}.
This logic implies that subglacial water which influences skin drag (related to form drag $\tau_b$) plays a less important role in ice dynamics when rougher, higher resolution bed elevation models (i.e. more form drag) are used, supporting studies where subglacial water inputs did not lead to significant speed ups in ice flow \citep[e.g.][]{SmithConnectedsubglaciallake2017}.

However, observed velocities over most of Pine Island Glacier are $>\SI{100}{\metre\per\year}$ (see Fig.~\ref{fig:topo_and_speed}b) which suggests a predominantly Coulomb-type (skin drag based) flow regime rather than a Weertman-type (form drag based) flow regime \citep[see Fig.~\ref{fig:1.3},][]{Minchewuniversalglacierslip2020}.
The decrease in basal drag $\tau_b$ (i.e. less skin drag) from using a higher resolution topography (i.e. more form drag) observed by \citet{Kyrke-SmithRelevanceDetailBasal2018} may be an artifact of using an unbounded Weertman-type sliding relation, as our experiments using a Schoof-type sliding law did not reproduce this behaviour.
In fact, Fig.~\ref{fig:basal_drag} shows that using a higher resolution DeepBedMap\_DEM over the lower resolution BedMachine (i.e. more form drag) actually led to a slight increase in inverted form drag $\tau_b$ (i.e. more skin drag), contrary to the findings of \citet{Kyrke-SmithRelevanceDetailBasal2018}.
The Weertman-style sliding law used by \citet{Kyrke-SmithRelevanceDetailBasal2018} in their inversion experiments lacks a parameterization of effective pressure $N$, thus excluding the possible effects of water cavitation \citep{BuddEmpiricalStudiesIce1979,GagliardiniFiniteelementmodelingsubglacial2007} or ice flow over deformable till \citep{Zoetsliplawglaciers2020}.
Their Weertman stress exponent value of $m=3$ also differs from that of \citet{JoughinRegularizedCoulombFriction2019} who recommends using $m=8$ for a Weertman slip relation to best match observed velocity data over Pine Island Glacier.
It thus remains to be seen whether our findings and those of \citet{Kyrke-SmithRelevanceDetailBasal2018} reporting on the effects of high resolution topography affecting skin drag can be generalized to more realistic sliding laws, and a wider area of the Antarctic Ice Sheet.

The inversion study here over the main trunk of Pine Island Glacier used a Schoof-type sliding relation on two beds with different spatial resolutions and roughness (i.e. different form drag).
Our findings suggest that skin drag (influenced by subglacial water and bed material properties) cannot be ruled out as an important factor in ice dynamics when using a high resolution bed, because basal drag ($\tau_b$, i.e. skin drag) did not noticeably decrease when more form drag was added with the use of a higher resolution (\SI{250}{\metre}) DeepBedMap\_DEM over the medium resolution (\SI{500}{\metre}) BedMachine (Fig.~\ref{fig:basal_drag}).
% \subsection{Future directions}
Future modelling work could include better parametrizations of effective pressure $N$, such as with an evolving subglacial hydrology model \citep[e.g.][]{SommersSHAKTISubglacialHydrology2018}.
Inversion experiments should also be carried out on even higher resolution ($<=\SI{100}{\metre}$) bed elevation models \citep[e.g.][]{Grahamhighresolutionsyntheticbed2017}.
A more thorough sensitivity analysis should be carried out to test different values of parameters such as $C_{\text{max}}$ in the Schoof sliding law, and also regularization parameters $\gamma_1$ and $\gamma_2$ (or others) to see how basal drag $\tau_b$ may change in magnitude or spatially.
A prognostic forward model run using these high resolution bed elevation models with a Schoof-type sliding law and Full Stokes (or higher order) stress balance equations would also be useful to assess the potential effects of different basal traction and slipperiness fields on Antarctic grounding line migration and its resulting contribution to sea level rise.


\chapter{Automated classification of active subglacial lakes in Antarctica from ICESat-2/ATLAS laser altimetry (2018--2020)}
\label{ch:4}
\chaptermark{ICESat-2 subglacial lakes}
% Active subglacial hydrology in Antarctica

\section{Subglacial Hydrology}

What is the behaviour of water underneath the ice, and what methods exist to map them.
Ultimately, how much of an influence does it actually have on Antarctic ice flow?

In Antarctica, water can be easily seen on the surface in some places, mostly close to dark, low-albedo areas like rock outcrops and blue ice regions \citep{KingslakeWidespreadmovementmeltwater2017}.
The bulk of liquid water however, lies hidden beneath the ice sheet.
Over 400 subglacial lakes have now been discovered in Antarctica, and we know from geomorphological evidence that water flows in subglacial channels underneath the ice sheet \citep{SiegertRecentadvancesunderstanding2016}.
Water in the cryospheric system is especially interesting primarily because of its fluid properties.
Compared to ice, water can flow a lot more quickly over short timescales, and ice that is in contact with water is more dynamic than it would otherwise be.

Observations of water under the ice sheet are usually done using remote sensing techniques, aided by the different transmission properties of ice, water and rock.
Surveying these features requires studying either elastic wave (e.g. seismic sounding) or electromagnetic wave (e.g. ice-penetrating radar) signals.
The waves may come from active or passive sources, and are detected from sensors deployed on the ground, in the air, or onboard of satellites in space.
We know that subglacial water can exist in three ways: 1) subglacial lakes, 2) subglacial channels, 3) subglacial aquifers \citep{ColleoniSpatiotemporalvariabilityprocesses2018}.
In 1967, the first direct documented evidence for a subglacial lake in Antarctica was obtained using airborne radio echo sounding data \citep{RobinInterpretationRadioEcho1969} from a joint programme between the UK Scott Polar Research Institute, the US National Science Foundation and the Technical University of Denmark, cumulating in the first subglacial lake inventory of 17 lakes \citep{OswaldLakesAntarcticIce1973}.
This was followed by a second inventory with 77 lakes \citep{SiegertinventoryAntarcticsubglacial1996}, and a third inventory with 145 lakes \citep{SiegertrevisedinventoryAntarctic2005}, all of which were detected using radio echo sounding techniques.

Since then, large (10+ km in diameter) `active' subglacial lakes have been detected based on vertical subglacial displacements using radar interferometry \citep{GrayEvidencesubglacialwater2005}, laser altimetry \citep{Smithinventoryactivesubglacial2009} and optical image differencing \citep{FrickerActiveSubglacialWater2007}, resulting in a fourth inventory that includes 379 subglacial lakes \citep{WrightfourthinventoryAntarctic2012}.
Subsequent discoveries have followed to bring the number to above 400 \citep[e.g.][]{WrightEvidencehydrologicalconnection2012,WrightSubglacialhydrologicalconnectivity2014,KimActivesubglaciallakes2016,RiveraSubglacialLakeCECs2015,SmithConnectedsubglaciallake2017}.

The way in which subglacial water routes itself beneath the ice is itself an important question.
Various types of subglacial pathways have been theorized over the years, ranging from fast flow in concentrated channels, to slower distributed flow spread out over a larger area (see Figure TODO).
The type of subglacial drainage system is likely to be influenced by the geology of the bedrock and the temperature profile of the ice column.
For soft bed types, water may incise into the rock to form semi-circular Nye channels or broad shaped canals, perhaps even flow along cavities, or if the rock is permeable, the water may well flow within the rock itself.
For hard bed types, water may incise instead into the ice bed, forming semi-circular Rothlisberger channels (TODO macron) or broad low channels, or perhaps flow as a thin sheet of water between the ice-rock interface.
These subglacial drainage structures are not static either, but are known to change between the two extremes of efficient and inefficient regimes over space and time that has important consequences for ice dynamics \citep{MullerVelocityfluctuationswater1973}.
However, it is important to keep in mind that the treatment of Antarctic glaciers or ice streams does differ from that of temperate glacier owing to the relative lack of input from surface meltwater.
The Antarctic subglacial water system is likely to be predominantly supplied from basal melt processes, such as from geothermal heat sources or perhaps from frictional heating under areas of fast ice flow.

Only in a handful of places has there been a comprehensive study that looked at the subglacial conditions of an Antarctic ice stream.
An area that was heavily focused on initially was the Whillans Ice Stream (formerly Ice Stream B) at the Siple Coast in West Antarctica where seismic surveys identified a water saturated, ~5 m thick porous till layer \citep{BlankenshipSeismicmeasurementsreveal1986} that could easily deform and explain the high surface velocities observed over that area \citep{AlleyDeformationtillice1986}.
Also in the vicinity at Subglacial Lake Whillans, ice penetrating radar \citep{ChristiansonSubglacialLakeWhillans2012} and active seismic \citep{HorganSubglacialLakeWhillans2012} surveys were conducted to constrain the thickness of the ice and water bodies as part of the Whillans Ice Stream Subglacial Access Research Drilling (WISSARD) project \citep{TulaczykWISSARDSubglacialLake2014}.

% TODO Debate about bed topography vs water
Indeed, while there is little question as to the importance of the bed, there has been considerable debate on the influence of water versus topographic controls on the flow of ice.
Water appears to be a clear interest in first decade of the 2000s.
On one hand, \citet{BellLargesubglaciallakes2007} found the onset of rapid flow at the downslope margins of four Recovery subglacial lakes...?

However, \citep{WinsborrowWhatcontrolslocation2010} says that topographic focusing is more important that meltwater/soft beds.

On one hand, \citet{SmithConnectedsubglaciallake2017} suggest that extra water had little or no influence on the speed of the lower trunk of Thwaites Glacier.

\citet{DiezBasalSettingsControl2018} ?Recovery/Slessor/Bailey Region!! Recovery Glacier topographically controlled in downstream area, upstream controlled by basal water?!
\citep{WrightSubglacialhydrologicalconnectivity2014,StearnsIncreasedflowspeed2008} Byrd glacier speedup from glacial lake drainage


% TODO how to detect?
% hydropotential formulation
% Estimate subglacial water routes https://www.mathworks.com/matlabcentral/fileexchange/55352-how-to-estimate-subglacial-water-routes
\subsection{Hydropotential}

Hydropotential (or hydraulic potential) refers to the static energy of water available at a particular time and place, and is also occassionally referred to in the literature as hydrostatic pressure.
Its calculation is basically a function of the amount of pressure exerted on a water body, located at a particular elevation relative to a reference datum.
By calculating hydropotential over a spatial surface, we can then derive the hydropotential gradients which provides us with a measure of the direction and tendency of water to flow if suitable conduits exist in its path.
Following the methods of \citet{ShreveMovementWaterGlaciers1972}, we calculate basal hydropotential at the ice-rock interface as follows:

\begin{equation}\label{eq:4.1}
  \phi = \phi_0 + p_w + \rho_wgz_b
\end{equation}

where the hydropotential at the base of the ice \gls{phi} is equal to an arbitrary constant $\phi_0$ plus pressure due to water $p_w$ plus the elevation potential term $\rho_wgz_b$.
The elevation potential term is made up of the density of water $\rho_w$ multiplied by the gravitational acceleration term $g$ multiplied by the bed elevation $z_b$.
If we assume no ice deformation, water pressure $p_w$ at the subglacial ice-rock interface can be approximated simply as the pressure induced by the overlying ice:

\begin{equation}\label{eq:4.2}
  p_w = \rho_i * g * (z_s - z_b)
\end{equation}

where the water pressure $p_w$ is equal to the density of ice $\rho_i$ multiplied by the gravitational acceleration term $g$ multiplied by the thickness of ice $z_s - z_b$ obtained from the ice surface elevation $z_s$ minus the ice bed elevation $z_b$.
Using a gravitational acceleration $g$ of $9.8ms^{-1}$, ice density $\rho_i$ of $917kgm^{-3}$, and water density $\rho_w$ of $1030kgm^{-3}$, we can substitute Equation \eqref{eq:4.2} into \eqref{eq:4.1} to obtain the following equation:

\begin{equation}\label{eq:4.3}
  \phi = 1107.4\left(\frac{917}{113} * z_s + z_b \right)
\end{equation}


% Started 18 July 2020
\subsection{Active Subglacial Lakes}

The volume of water in subglacial lakes can change rapidly over short periods of times, happenning in a matter of days or weeks, and we call these active subglacial lakes.
As these dynamic changes in water volume typically results in changes in the ice elevation surface, we can detect active subglacial lakes more easily using satellite-based techniques, in contrast to passive subglacial lakes that can only be detected via ground-based radar and seismic surveys.


\subsubsection{Measuring Elevation (h) over the Antactic Ice Sheet}

Several satellite sensors exist to measure ice surface height, such as laser and radar altimeters, or optical photogrammetry.
A comparison of 3 main methods of satellite-based is given below:

\begin{enumerate}
  \item Laser altimeter:
  \begin{itemize}
    \item Most direct measurement of ice surface, no firn penetration
    \item Uses green 532nm wavelength light
    \item Measurements may need to be corrected by slope
    \item Affected by cloud cover
    \item E.g. ICESat-2 reaches 88°S, ICESat-1 reaches 86°S
  \end{itemize}

  \item Radar altimeter:
  \begin{itemize}
    \item Direct measurement of ice surface, but need to be corrected for firn penetration
    \item Uses radar (X band)
    \item Unaffected by cloud cover
    \item E.g. Cryosat-2 reaches 88°S
  \end{itemize}

  \item Photogrammetry:
  \begin{itemize}
    \item Somewhat direct measurement of ice surface height
    \item Requires multiple passes over the same location to derive the parallax measurement
    \item Uses optical sensors (e.g Red, Green, Blue, Infrared bands)
    \item Requires high radiometric resolution over ice sheets that are mostly white
    \item Affected by cloud cover
    \item High temporal resolution compared to altimeters
    \item E.g. Landsat-8 reaches reaches 8?°S TODO
  \end{itemize}
\end{enumerate}

We will focus on laser altimeters, specifically the ICESat-2/ATLAS satellite sensor, as it is the most direct method of obtaining ice surface height information.
Satellite altimeters work by beaming an electromagnetic wave pulse down to the Earth surface, where it is reflected off the ice surface and goes back to the satellite sensor.
The time it takes from when the beam is sent out to when it is recorded again by the sensor allows us to measure the distance between the satellite and the ice surface:

\begin{equation}\label{eq:4.4}
  d = \frac{t * c}{2}
\end{equation}

where multiplying the speed of light $c$ by the time \gls{t} taken, and dividing by two, gives the distance $d$ from the satellite to the ground.
The elevation of the ice surface is then given by:

\begin{equation}\label{eq:4.5}
  h = z_p - d
\end{equation}

where the ice surface elevation $h$ is given by the the satellite platform's elevation $z_p$ minus the distance from the satellite platform to the ice surface $d$.

% TODO insert illustration of satellite altimeter measurement

\subsubsection{Ice Surface Elevation Change over Time (dh)}

Ice surface elevation $h$ changes over time \gls{t} due to several processes, both on the surface of the ice, and what happens underneath at the bed.
It is measured as follows:

\begin{equation}\label{eq:4.6}
  dh = h_1 - h_0
\end{equation}

where the change in ice surface elevation $dh$ between two points in time is the difference between heights $h_1$ and $h_0$ respectively over time $t=1$ and $t=0$.
This ice surface elevation change can be broken down into two components:

\begin{equation}\label{eq:4.7}
  dh = dh_s + dh_b
\end{equation}

where the total change in ice surface elevation $dh$ is equal to the change in elevation due to surface processes $dh_s$ plus the change in elevation due to basal proccesses $dh_b$.

Our study focuses on the subglacial component, but because we are only able to directly measure ice surface changes $dh$ using satellites, it is important to recognize the different processes that lead to the total change in ice surface elevation over time.
Specifically, the magnitude of vertical elevation change, and the timespan in which this elevation change occurs are quite different between surface and basal processes.
We can thus use these different patterns in the elevation timeseries data to determine the cause of elevation changes.

On the surface (air-ice interface), snow accummulation can lead to an increase in elevation, while a decrease in elevation can come from snow compaction or mass wasting processes.
Surface proccesses are seasonally dependent, with an increase in elevation due to snow accummulation over the austral winter, and a decrease in elevation due to snow compaction in the austral summer.
In Antarctica, accummulation tends to be a slow and gradual process due the low precipitation over the continent, in the order of ~2 mm a year (TODO cite Arthern).
Mass wasting processes typically occur in fast flowing ice regions like glaciers and ice streams.
Over a year, the net contribution of these surface processes are typically in the order of centimetres to tens of centimetres a year.
Over a decade, these may amount to about a metre or less of elevation change, and we can this stable linear trend in most parts of the Antarctic continent over the satellite era (TODO cite 2020 Science paper on ice loss).

% TODO insert figure of some slowly accummulating/mass wasting region in Antarctica with time on x-axis and elevation on y-axis

On the bottom (ice-rock interface), changes in the volume of subglacial water, such as due to subglacial lake drainage or filling events, can also lead to ice surface elevation changes.
The mechanism of subglacial lake drainage and filling events are known (TODO or is it?, find citation), though predicting these events remain an elusive task.
The drainage or filling however, can take place rapidly at time scales of a few days or weeks, with ice surface elevation changes up to a few metres observed.
The pattern of sudden elevation change over a short timescale caused by subglacial water volume changes typically occurs over a limited geographical region of a few square kilometres (?TODO fact check).
This stands in contrast with the background gradual trend in elevation change caused by surface processes.
It is this anomalous pattern which is used as the basis of active subglacial lake detection.

\subsubsection{Rate of Elevation Change over Time (dh/dt)}

In order to reliably detect subglacial lake drainage or filling events, we need a metric that is robust enough to work on a given dataset of time series elevation measurements.
Ideally, the metric would be able to capture both the magnitude of vertical elevation change, and the timespan in which the elevation change has occured.
We will present and discuss two options of varying computational complexity to do this.

The magnitude of vertical elevation change can be represented directly in the form of a height range:

\begin{equation}\label{eq:4.8}
  h_{range} = h_{max} - h_{min}
\end{equation}

where the height range $h_{range}$ is equal to the maximum height $h_{max}$ minus the minimum height $h_{min}$ at a given point.
Computationally, this has $O(n)$ complexity (TODO double-check $O(n)$ notation).

However, there are a couple of disadvantages to this approach:

\begin{itemize}
  \item It is prone to outliers. For example, an anomalously high elevation caused by cloud interference can lead to an incorrect $h_{range}$ value.
  \item No sense of direction. There is no way to tell whether the $h_{range}$ value represents a surface rising or lowering.
\end{itemize}

That said, the calculation is fast to compute and scales linearly with the number of data points.
It can be used to quickly locate areas that have experienced a significant amount of elevation change.
Areas with high $h_{range}$ values can then be analyzed further using the linear regression method we describe next.

The rate of ice surface elevation change over time can calculated using an ordinary least squares formula:

\begin{equation}\label{eq:4.9}
  \frac{dh}{dt} = OLS( h_0, t_0, h_1, t_1, \dots, h_n, t_n )
\end{equation}

where the rate of elevation change over time $\frac{dh}{dt}$ is found through a linear regression of data points elevation $h_0$ at time $t_0$ until $h_n$ at time $t_n$.
The slope of the line of best fit represents $\frac{dh}{dt}$, and is typically given in metres per year.
Computationally, this has $O(\log n)$ complexity (TODO double-check $O(\log n)$ notation).
This makes it more computationally expensive than the $h_{range}$ calculation, and thus more effort is require to scale it to continent-wide analysis.

% TODO insert figure of linear regression of one ATL11 point

The benefits of this approach are:

\begin{itemize}
  \item The direction in which the ice surface elevation is changing is preserved, as we can tell whether the elevation is going up or down.
  \item More robust to outliers than $h_{range}$.
\end{itemize}

However, care should be taken when interpreting rate of elevation change $\frac{dh}{dt}$ for any given point.
Fine scale details in the elevation change may be masked, for example, if the surface were to rise at $t_1$ and fall at $t_2$, the $\frac{dh}{dt}$ value may not indicate a change.
Care should be taken to ensure the $\frac{dh}{dt}$ calculation starts and ends at about the same time in the year to account for seasonal effects.
While this method is less sensitive to outliers, it is still a good idea to remove points with anomalous height values before calculating $\frac{dh}{dt}$.
This can be done by filtering out the points based on quality assessment criteria such as the photon return signal.
Alternatively, a weighted linear regression can be used to give more weight to good data points (i.e. those less affected by clouds), and less weight to the anomalous points.

\subsubsection{Detecting Active Subglacial Lakes}

The algorithm we use to detect active subglacial lakes builds upon that of those used in the ICESat-1 era \citep{FrickerActiveSubglacialWater2007,Smithinventoryactivesubglacial2009}, updated for the ICESat-2 era.
In terms of active subglacial lake detection, the areas of improvements of the current generation ICESat-2 satellite include:

\begin{itemize}
  \item the six laser beam setup allow us to better handle anomalous heights due to cross-track slopes that was a major issue in ICESat-1.
  \item there is an order of magnitude increase in track density, and thus greater coverage of Northern areas of the Antarctica continent, notably near the coastal areas.
  \item the coverage hole over the South Pole has narrowed, with an additional 2 degrees of latitude from 86° S to 88° S now covered by ICESat-2
\end{itemize}

These new features in ICESat-2 allow us to simplify the active subglacial lake detection algorithm in \citep{Smithinventoryactivesubglacial2009} considerably in terms of less cross-track slope correction.
However, the increased data density of ICESat-2 over ICESat-1 also requires a more deliberate data engineering process.
While we continue to filter out data that are problematic due to clouds, there is now a specific intermediate step to filter out areas that are less likely to harbour active subglacial lakes.

Our algorithm is as follows:

\begin{enumerate}
  \item For an ATL11 point, calculate $h_{range}$. Keep ATL11 points that have $h_{range} > 0.5 m$, drop the rest
  \item For the remaining ATL11 points, calculate $\frac{dh}{dt}$. Keep ATL11 points that have $\frac{dh}{dt} > 0.5 m y^{-1}$, drop the rest
  \item For the remaining ATL11 points, check that there exists a spatial cluster of points with similar $\frac{dh}{dt}$ values surrounding it (1 sq km? TODO). Keep ATL11 points that are within high $\frac{dh}{dt}$ spatial clusers, drop the rest.
  \item For the remaining ATL11 points, cross correlate the areas with additional information such as ice velocity data, hydropotential calculations, image differencing, etc.
  \item Generate polygons of active subglacial lake boundaries, with information on estimated ice volume gained/lost over \gls{t} amount of time. TODO properly.
\end{enumerate}

Owing to the dynamic nature of these subglacial lake activity, wherein lakes can drain in a matter of weeks or months, we have automated the procedure to make it easier to get into more detailed analysis sooner.
The automation allows us to detect such events every season, or every month, depending on data availability.

One thing that the increased ICESat-2 spatial data density has shown us over ICESat-1 is that the distribution of active subglacial lakes over Antarctica is greater than once thought.
Instead of ~145 subglacial lakes with a combined volume of ~?? cubic metres \citep{Smithinventoryactivesubglacial2009}, we have detected ~??? subglacial lakes with a combined volume of ~?? cubic metres.
This has implications for (to be continued).


\chapter{Conclusion}
% conclusions

This thesis maps out the subglacial world of Antarctica using satellite measured surface observations.
A series of machine learning, inverse modelling and remote sensing methods were presented, utilizing high spatial resolution ($<= \SI{250}{\metre}$) datasets to investigate Antarctica's subglacial topography and hydrology, with implications for ice flow dynamics and future sea level projections.
This section synthesizes the discoveries made in light of answering our original research questions, highlights gaps remaining in both our observations and numerical models of subglacial Antarctica, and presents a path forward for future scholars.

\section{Research Questions}

\begin{enumerate}
  \item Can we use a super-resolution convolutional neural network to create a higher spatial resolution (\SI{250}{\metre}) bed elevation map of Antarctica?

  A deep neural network called DeepBedMap was designed to super-resolve the bed elevation of Antarctica \citep{LeongDeepBedMapdeepneural2020}, with the model architecture adapted from an Enhanced Super-Resolution Generative Adversarial Network \citep{WangESRGANEnhancedSuperResolution2019}.
  The DeepBedMap model is trained on areas in Antarctica where high-resolution (\SI{250}{\metre}) ground-truth bed elevation grid images are available, and then applied across the entire Antarctic continent.
  Inputs to the model are a prior low-resolution (\SI{1}{\kilo\metre}) BEDMAP2 image, and conditional remote sensing images of ice surface elevation, velocity and snow accumulation.
  Model training occurs via an iterative error minimization approach with a custom perceptual loss function that pushes the predicted images to be as close as possible to the ground-truth images.
  The model output is a four-times-upsampled super-resolution (\SI{250}{\metre}) bed elevation model of Antarctica called DeepBedMap\_DEM that preserves detail in bed roughness and can be used for catchment- to continent-scale ice sheet modelling studies.

  \item What effect does a rough surface, high-resolution (\SI{250}{\metre}) bed topography have on the friction parameters of an ice sheet model?

  To examine the relative contributions of form drag and skin drag which opposes the driving stress of a glacier, a basal inversion experiment was performed on two different bed topographies - a rougher high-resolution (\SI{250}{\metre}) DeepBedMap\_DEM \citep{LeongDeepBedMapdeepneural2020} and a smoother medium-resolution (\SI{500}{\metre}) BedMachine Antarctica \citep{MorlighemDeepglacialtroughs2019}.
  The inversion was conducted via an iterative least-squares control method using the Ice-sheet and Sea-level System Model \citep{LarourContinentalscalehigh2012} with a Coulomb-limited Schoof-type sliding law.
  The inverted basal fields between DeepBedMap\_DEM and BedMachine appear to be within the one sigma error margin, though both basal drag and slipperiness fields are slightly higher while effective pressure is generally lower for the inversions ran on the higher-resolution DeepBedMap\_DEM.
  The results here do not support the findings of previous studies using a Weertman-style sliding relation by \citet{Kyrke-SmithRelevanceDetailBasal2018} who indicated that skin drag was reduced when high resolution bed topography was incorporated into the model.
  Our findings in this chapter highlight that more research is needed over diverse subglacial settings in West Antarctica to investigate the role of subglacial hydrology which influences skin drag, on high resolution bed topographies ($\leq\SI{100}{\metre}$ which influences form drag.

  \item Where does water drain and accumulate underneath the Antarctic Ice Sheet, how much volume exists, and at what timescales do these processes occur?

  A map of active subglacial lakes in Antarctica is presented, detected using ICESat-2 laser altimetry data over the 2018-2020 time period.
  Active subglacial lake locations were determined via an unsupervised density-based classification method on ICESat-2 point cloud data pre-processed to highlight anomalous ice surface elevation change rates.
  The algorithm yielded a total of 149 active subglacial lakes, including 12 new lakes in the 86--88°S area not detected by the previous ICESat (2003-2009) mission.
  We detailed a cascading pattern of drain-fill activity over the Whillans Ice Stream central basin at the Siple Coast, and revealed multi-lobe subglacial lake clusters separated by ridges using the high resolution ($<\SI{40}{\metre}$ along-track spacing) ICESat-2 laser altimetry data.
  An interesting observation was the unusually high rate of surface elevation uplift and slow descent at Subglacial Lake Whillans IX ($\sim\SI{7}{\metre}$ vertical rise over 3 months) and Whillans 7 ($\sim\SI{-6}{\metre}$ lowering over 11 months) which stands in contrast with previous studies over the Whillans Ice Stream catchment area that tend to show slow fill and rapid drainage \citep{SiegfriedThirteenyearssubglacial2018,SiegfriedEpisodicicevelocity2016}.

\end{enumerate}

\section{Future work}

Advances in glaciology are driven by the availability of newer, high quality datasets combined with better, physical numerical models.
In the following, we highlight data gaps to be filled, and missing components in our understanding of the physical mechanisms of ice flow over the subglacial terrain of Antarctica.

\subsection{Towards BEDMAP3 - More data to super-resolve the bed topography of Antarctica}

In Chapter \ref{ch:2}, we indicated that there is missing coverage of bed elevation data in parts of Antarctica, and that Radio-echo sounding (\gls{RES}) is the best tool available to fill in these gaps.
Swath processing of RES data \citep[e.g.][]{HolschuhLinkingpostglaciallandscapes2020} should become a priority, and new acquisitions should target a diverse range of bed and flow types.
Furthermore, we should continue to survey formerly glaciated beds around the margins of Antarctica using ship-based swath bathymetry instruments and on land in areas like the former Laurentide ice sheet using LIDAR.
Continued measurements by remote sensing satellites is also needed to patch up data gaps in surface elevation, velocity and snow accumulation.

\subsection{Coupling ice flow models with evolving subglacial hydrological models}

In Chapter \ref{ch:3}, our basal inversion experiment suggested that subglacial water which influences skin drag cannot be ruled out as an important factor that affects ice dynamics when using a high resolution bed.
Hence, future modelling studies could include a coupling with an evolving subglacial hydrology model \citep[e.g.][]{SommersSHAKTISubglacialHydrology2018} to better parametrize the effective pressure field in a Coulomb-limited sliding law.
More prognostic forward model runs should be done on high resolution bed elevation models with a Schoof-type sliding law and Full Stokes (or higher order) stress balance equations to assess the potential effects of different basal traction and slipperiness fields on Antarctic grounding line migration and its resulting contribution to sea level rise.

\subsection{Continuous active subglacial lake time-series data}

In Chapter \ref{ch:4}, we used the ICESat-2 laser altimeter to create an inventory of active subglacial lakes for a short time period from 2018-2020, but future work should extend the record both backwards in time to the ICESat era (2003-2009) and into the future such as with the CRISTAL radar altimeter mission \citep{KernCopernicusPolarIce2020}.
A priority for the remote sensing community will be to mitigate against any interruptions in the data record (such as between ICESat and Cryosat-2 previously), utilizing oversnow geophysical surveys or continuous GPS stations if necessary on key sites to minimize loss of continuous coverage.
One critical piece of work will lie in calibrating differences in elevation measurements across laser- and radar-based satellite altimeters \citep[c.f.][]{SiegfriedThirteenyearssubglacial2018}, for the currently operational ICESat-2 and Cryosat-2 \citep{BruntComparisonsSatelliteAirborne2020} and in the future as new satellite altimeters come online.

\section{Concluding remark}

The objective of this thesis were to 1) create a high-resolution (\SI{250}{\metre}) map of Antarctica's bed, 2) investigate how high-resolution bed topography affects basal friction, and 3) map out the active subglacial hydrological system of Antarctica.
These are addressed using a series of novel machine learning and automated satellite remote sensing methods that use surface observations to infer the nature of Antarctica's subglacial topography and hydrology.


\appendix
% Appendix

\chapter{Details of loss function components}
\chaptermark{Loss function}
\label{appendix:A}

The loss function, or cost function, is a mathematical function that maps a set of input variables to an output loss value.
The loss value can be thought of as a weighted sum of several error metrics between the neural network's prediction and the expected output or ground truth.
It is this loss value which we want to minimize so as to train the neural network model to perform better, and we do this by iteratively optimizing the parameters in the loss function.
Following this are the details of the various loss functions that make up the total loss function of the DeepBedMap generative adversarial network model.

\section{Content Loss}

To bring the pixel values of the generated images closer to those of the ground truth, we first define the content-loss function~$L_1$.
Following ESRGAN \citep{WangESRGANEnhancedSuperResolution2019}, we have

\begin{equation}\label{eq:A1}
  L_1 = \dfrac{1}{n} \sum\limits_{i=1}^n ||\hat{y}_i - y_i||_{1}~,
\end{equation}

where we take the mean absolute error between the generator network's predicted value~$\hat{y}_i$ and the ground-truth value~$y_i$, respectively, over every pixel $i$.

\section{Adversarial Loss}

Next, we define an adversarial loss to encourage the production of high-resolution images~$\hat{y}$ closer to the manifold of natural-looking digital-elevation-model images.
To do so, we introduce the standard discriminator in the form of $D(y) = \sigma(C(y))$, where~$\sigma$ is the sigmoid activation function and $C(y)$ is the raw, non-transformed output from a discriminator neural network acting on high-resolution image~$y$.
The ESRGAN model \citep{WangESRGANEnhancedSuperResolution2019}, however, employs an improved relativistic-average discriminator \citep{Jolicoeur-Martineaurelativisticdiscriminatorkey2018} denoted by~$D_{\text{Ra}}$.
It is defined as $D_{\text{Ra}}(y,\hat{y}) = \sigma(C(y) - \mathbb{E}_{\hat{y}}[C(\hat{y})])$, where $\mathbb{E}_{\hat{y}}[\cdot]$ is the arithmetic mean operation carried out over every generated image $\hat{y}$ in a mini batch.
We use a binary cross-entropy loss as the discriminator's loss function defined as follows:

\begin{equation}\label{eq:A2}
  L_{\mathrm{D}}^{\text{Ra}} = - \mathbb{E}_y[\ln(D(y,\hat{y}))] - \mathbb{E}_{\hat{y}}[\ln(1 - D(\hat{y},y))].
\end{equation}

The generator network's adversarial loss is in a symmetrical form:

\begin{equation}\label{eq:A3}
  L_{\mathrm{G}}^{\text{Ra}} = - \mathbb{E}_y[\ln(1 - D(y,\hat{y}))] - \mathbb{E}_{\hat{y}}[\ln(D(\hat{y},y))].
\end{equation}

\section{Topographic Loss}

We further define a topographic loss so that the elevation values in the super-resolved image make topographic sense with respect to the original low-resolution image.
Specifically, we want the mean value of each 4\,$\times$\,4 grid on the predicted super-resolution (DeepBedMap) image to closely match its spatially corresponding \SI{1}{pixel}\,$\times$\,\SI{1}{pixel} area on the low-resolution (BEDMAP2) image.

First, we apply a 4\,$\times$\,4 mean pooling operation on the generator network's predicted super-resolution image:

\begin{equation}\label{eq:A4}
  \bar{\hat{y}}_j = \dfrac{1}{n} \sum\limits_{i=1}^n \hat{y}_i~,
\end{equation}

where $\bar{\hat{y}}_j$ is the mean of all predicted values~$\hat{y}_i$ across the 16 super-resolved pixels~$i$ within a 4\,$\times$\,4 grid corresponding to the spatial location of 1 low-resolution pixel at position $j$.
Following this, we can compute the topographic loss as follows:

\begin{equation}\label{eq:A5}
  L_{\mathrm{T}} = \dfrac{1}{m} \sum\limits_{i=1}^m ||\bar{\hat{y}}_j - x_j||_{1}~,
\end{equation}

where we take the mean absolute error between the mean of the 4\,$\times$\,4 super-resolved pixels calculated in Eq.~(\ref{eq:A4}) $\bar{\hat{y}}_j$ and those of the spatially corresponding low-resolution pixel~$x_j$, respectively, over every low-resolution pixel~$j$.

\section{Structural Loss}

Lastly, we define a structural loss that takes into account luminance, contrast and structural information between the predicted and ground-truth images.
This is based on the structural similarity index \citep[SSIM;][]{WangImageQualityAssessment2004} and is calculated over a single window patch as

\begin{equation}\label{eq:A6}
  \text{SSIM} (\hat{y}, y) = \dfrac{(2\mu_{\hat{y}}\mu_y + c_1)(2\sigma_{{\hat{y}}y} + c_2)}{(\mu_{\hat{y}}^2 + \mu_y^2 + c_1)(\sigma_{\hat{y}}^2 + \sigma_y^2 + c_2)},
\end{equation}

where $\mu_{\hat{y}}$ and $\mu_y$ are the arithmetic mean of predicted image~${\hat{y}}$ and ground-truth image~$y$, respectively, over a single window that we set to 9~pixels\,$\times$\,9~pixels; $\sigma_{{\hat{y}}y}$ is the covariance of~${\hat{y}}$ and~$y$; $\sigma_{\hat{y}}^2$ and $\sigma_y^2$ are the variances of~${\hat{y}}$ and~$y$, respectively; and~$c_1$ and~$c_2$ are two variables set to $0.01^2$ and $0.03^2$ to stabilize division with a weak denominator.
Thus, we can formulate the structural loss as follows:

\begin{equation}\label{eq:A7}
  L_{\mathrm{S}} = 1 - \dfrac{1}{p} \sum\limits_{i=1}^p \text{SSIM} (\hat{y}, y)_p~,
\end{equation}

where we take $1$ minus the mean of all structural similarity values $\text{SSIM}(\hat{y}, y)$ calculated over every patch~$p$ obtained via a sliding window over the predicted image~${\hat{y}}$ and ground-truth image~$y$.

\section{Total Loss Function}

Finally, we compile the loss functions for the discriminator and generator networks as follows:

\begin{align}
  & L_{\mathrm{D}} = L_{\mathrm{D}}^{\text{Ra}}, \label{eq:A8}\\
  & L_{\mathrm{G}} = \eta L_1 + \lambda L_{\mathrm{G}}^{\text{Ra}} + \theta L_{\mathrm{T}} + \zeta L_{\mathrm{S}}~, \label{eq:A9}
\end{align}

where $\eta$, $\lambda$, $\theta$ and $\zeta$ are the scaled weights for the content $L_1$, adversarial $L_{\mathrm{D}}$, topographic $L_{\mathrm{T}}$ and structural losses $L_{\mathrm{S}}$, respectively (see Table~\ref{table:B1} for values used).
The loss functions $L_{\mathrm{D}}$ and $L_{\mathrm{G}}$ are minimized in an alternate $1:1$ manner so as to solve the entire generative adversarial network's objective function defined in Eq.~(\ref{eq:4}).


\chapter{Neural Network Training Details} \label{appendix:B}

The neural networks were developed using Chainer v7.0.0 \citep{TokuiChainerDeepLearning2019} and trained using full-precision (floating point 32) arithmetic.
Experiments were carried out on four graphical processing units (GPUs), specifically two Tesla P100 GPUs and two Tesla V100 GPUs.
On the Tesla V100 GPU setup, one training run with about 150 epochs takes about 30\,min.
This is using a batch size of 128 on a total of 3826 training image tiles, with 202 tiles reserved for validation, i.e. a $95/5$ training/validation split.
We next describe the method used to evaluate each DeepBedMap candidate model, as well as the high-level way in which we semi-automatically arrived at a good model via semi-automatic hyperparameter tuning.

%TB1
\begin{table}[b!]
  \small
  \caption{Optimized hyperparameter settings.}
  \label{table:B1}
  \begin{tabularx}{\textwidth}{lrr}
  \hline
  Hyperparameter                                       & Setting             & Tuning range                           \\
  \hline
  Learning rate (for both generator and discriminator) & $1.7\times 10^{-4}$ & $2\times 10^{-4}$ to $1\times 10^{-4}$ \\
  Number of residual-in-residual blocks                & 12                  & 8 to 14                                \\
  Mini-batch size                                      & 128                 & 64 or 128                              \\
  Number of epochs                                     & 140                 & 90 to 150                              \\
  Residual scaling                                     & 0.2                 & 0.1 to 0.5                             \\
  Content-loss weighting $\eta$                        & $1\times 10^{-2}$   & Fixed                                  \\
  Adversarial-loss weighting $\lambda$                 & $2\times 10^{-2}$   & Fixed                                  \\
  Topographic-loss weighting $\theta$                  & $2\times 10^{-3}$   & Fixed                                  \\
  Structural-loss weighting $\zeta$                    & 5.25                & Fixed                                  \\
  He normal initialization scaling                     & 0.1                 & Fixed                                  \\
  Adam optimizer epsilon                               & 0.1                 & Fixed                                  \\
  Adam optimizer beta1                                 & 0.9                 & Fixed                                  \\
  Adam optimizer beta2                                 & 0.99                & Fixed                                  \\
  \hline
  \end{tabularx}
\end{table}

To check for overfitting, we evaluate the generative adversarial network model using the validation dataset after each epoch using two performance metrics~-- a peak signal-to-noise ratio (PSNR) metric for the generator and an accuracy metric for the discriminator.
Training stops when these validation performance metrics show little improvement, roughly at 140~epochs.

Next, we conduct a full evaluation on an independent test dataset, comparing the model's predicted grid output with actual ground-truth \textit{xyz} points.
Using the ``grdtrack'' function in Generic Mapping Tools v6.0 \citep{WesselGenericMappingTools2019}, we obtain the grid elevation at each ground-truth point and use it to calculate the elevation error on a point-to-point basis.
All of these elevation errors are then used to compute a root mean square error (RMSE) statistic over this independent test site.
This RMSE value is used to judge the model's performance in relation to baseline bicubic interpolation and is also the metric minimized by a hyperparameter optimization algorithm which we will describe next.

Neural networks contain a lot of hyperparameter settings that need to be decided upon, and generative adversarial networks are particularly sensitive to different hyperparameter settings.
To stabilize model training and obtain better performance, we tune the hyperparameters (see Table~\ref{table:B1}) using a Bayesian approach.
Specifically, we employ the Tree-structured Parzen Estimator \citep{BergstraAlgorithmsHyperparameterOptimization2011} from the Optuna v2.0.0 \citep{AkibaOptunaNextgenerationHyperparameter2019} library with default settings as per the Hyperopt library \citep{BergstraHyperoptPythonlibrary2015}.
Given that we have four GPUs, we choose to parallelize the hyperparameter tuning experiments asynchronously between all four devices.
The estimator first conducts 20 random experimental trials to scan the hyperparameter space, gradually narrowing down its range to a few candidate hyperparameters in subsequent experiments.
We set each GPU to run a target of 60 experimental trials (i.e. a total of 240), though unpromising trials that have exploding or vanishing gradients are pruned prematurely using the Hyperband algorithm \citep{LiHyperbandNovelBanditBased2018} to save on time and computational resources.
The top models from these experiments undergo further visual evaluation, and we continue to conduct further experiments until a suitable candidate model is found.


\chapter{Hydropotential} \label{sec:hydropotential}

% hydropotential formulation
% Estimate subglacial water routes https://www.mathworks.com/matlabcentral/fileexchange/55352-how-to-estimate-subglacial-water-routes

Hydropotential (or hydrostatic pressure) refers to the static energy of water available at a particular time and place.
It is a function of the amount of pressure exerted on a water body, located at a particular elevation relative to a reference datum.
By calculating hydropotential over a spatial surface, we can then derive the hydropotential gradients which provides us with a measure of the direction and tendency of water to flow if suitable conduits exist in its path.
Following the methods of \citet{ShreveMovementWaterGlaciers1972}, basal hydropotential \gls{phi} is calculated as follows:

\begin{equation}\label{eq:4.11}
  \phi = \phi_0 + p_w + \rho_wgz_b
\end{equation}

where \gls{phi} denotes hydropotential at the base of the ice, $\phi_0$ is an arbitrary constant, $p_w$ is water pressure and $\rho_wgz_b$ is the elevation potential term.
The elevation potential term is made up of the density of water $\rho_w$ multiplied by the gravitational acceleration term $g$ multiplied by the bed elevation $z_b$.
In most cases, subglacial water pressure $p_w$ can be approximated as the pressure induced by the overlying ice (overburden pressure):

\begin{equation}\label{eq:4.12}
  p_w = \rho_i * g * (z_s - z_b)
\end{equation}

where the static water pressure $p_w$ is equal to the density of ice $\rho_i$ multiplied by the gravitational acceleration term $g$ multiplied by the thickness of ice $z_s - z_b$ obtained from the ice surface elevation $z_s$ minus the ice bed elevation $z_b$.
Using a gravitational acceleration $g$ of \SI{9.8}{\metre\per\second}, ice density $\rho_i$ of \SI{917}{\kilo\gram\per\metre\cubed}, and water density $\rho_w$ of \SI{1000}{\kilo\gram\per\metre\cubed}, we can substitute Equation \eqref{eq:4.12} into \eqref{eq:4.11} to obtain the following equation:

\begin{equation}\label{eq:4.13}
  \phi = 813.4\left(\frac{917}{83} * z_s + z_b \right)
\end{equation}

where the ice surface elevation $z_s$ is about $\frac{917}{83} = 11.05$ times more important than bed surface elevation $z_b$ for its effect on hydropotential $\phi$.
This is an oft-cited constant, and varies across the literature from as low as $8$ to as high as $11$ depending on what densities of water $\rho_w$ and ice $\rho_i$ are used in the calculation.

% See CarterAntarcticsubglaciallakes2017 Section 3.1, very good hydropotential equation description

% Water under the ice sheet can be detected by remote sensing techniques.
% This relies on the different transmission properties of ice, water and rock.
% Surveying these features requires studying either elastic wave (e.g. seismic sounding) or electromagnetic wave (e.g. ice-penetrating radar) signals.
% The waves may come from active or passive sources, and are detected from sensors deployed on the ground, in the air, or onboard of satellites in space.


\printbibliography

\end{document}
